\ifdefined\submissionmode
  \documentclass[acmsmall,screen,review,anonymous]{acmart}
\else
  \documentclass[acmsmall,screen,nonacm]{acmart}
\fi

\usepackage{tabularx}
\usepackage{bbm}
\usepackage{bbold}
\usepackage{stmaryrd}
\usepackage{booktabs}
\usepackage{enumitem}
\usepackage{mathpartir}
\usepackage{subcaption}
\usepackage{tikz-cd}
\usepackage{xspace}
% \finalmode (set by main-final.tex) renders the clean version.
\ifdefined\finalmode
  \usepackage[final,commandnameprefix=ifneeded,commentmarkup=footnote]{changes}
\else
  \usepackage[commandnameprefix=ifneeded,commentmarkup=footnote]{changes}
\fi
\usepackage{changebar}
\ifdefined\finalmode\nochangebars\fi

%\geometry{paperwidth=8.3in, paperheight=11.7in} % force to A4 for now
\settopmatter{printacmref=false}
\citestyle{acmauthoryear}
\raggedbottom

\input{macros}

\begin{document}

\title{Data Provenance as Automatic Differentiation}

\author{Robert Atkey}
\email{robert.atkey@strath.ac.uk}
\orcid{0000-0002-4414-5047}
\affiliation{%
  \institution{University of Strathclyde}
  \city{Glasgow}
  \country{UK}}

\author{Roly Perera}
\email{roly.perera@cl.cam.ac.uk}
\orcid{0000-0001-9249-9862}
\affiliation{%
  \institution{University of Cambridge}
  \city{Cambridge}
  \country{UK}
}
\additionalaffiliation{%
   \institution{University of Bristol}
   \city{Bristol}
   \country{UK}
}

\begin{abstract}
  Automatic differentiation (AD) computes the derivative of a program alongside the program itself,
  as a linear map between tangent spaces, propagated forwards or backwards along an execution. We present a
  semantic framework that models \emph{data provenance} via the same construction: taking scalars from a
  commutative semiring of dependency information rather than the real numbers, the derivative of a program
  becomes a linear map between spaces of approximations of its input and output. The choice of semiring
  determines the notion of provenance. Over the two-element Boolean algebra, the Jacobian of a program
  records which input positions each output position may depend on, and composing Jacobians forwards or
  backwards is \emph{dependency analysis} in the manner of forward- and reverse-mode AD. More generally, over
  distributive lattices the Jacobian and its transpose propagate dependency information forwards and
  backwards as a \emph{conjugate pair} of maps; when the lattice is a Boolean algebra, the two directions are
  moreover related by adjunction, recovering an approach called \emph{\GPS}. We interpret a higher-order total
  functional language in this framework, prove that every program of first-order type denotes such a Jacobian,
  and instantiate the semiring to obtain dependency tracking (Booleans), automatic differentiation (reals),
  and quantitative interval provenance (the tropical semiring) as examples. All results are formalised in
  Agda.
\end{abstract}
\maketitle

\section{Introduction}
\label{sec:introduction}

To audit any computational process, we need robust and well-founded notions of \emph{provenance} to track how
data are used. This allows us to answer questions like ``Where did these data come from?'', ``Why are these
data in the output?'' and ``How were these data computed?''. Provenance tracking has a wide range of
applications, from debugging and program comprehension~\cite{buneman95,cheney07} to improving reproducibility
and transparency in scientific workflows~\cite{kontogiannis08}. \emph{Program slicing}, first proposed
by~\citet{weiser81}, is a collection of techniques for provenance tracking that attempts to take a run of a
program and areas of interest in the output, and turn them into the subset of the input and the program that
were responsible for generating those specific outputs.

Underlying all of these questions is a more basic one: how does the output of a program respond when
its input changes? For programs that compute with real numbers, calculus provides a canonical answer: the
derivative, a linear map between tangent spaces describing how the output varies as the input varies, with
\emph{automatic differentiation} (AD) being a computational technique for computing such derivatives alongside
the program itself~\cite{siskind08,elliott18,vakar22}. For data provenance and dependency analysis, where the
inputs are database tuples or elements of a data structure, the question is often framed in a more qualitative
way: if we perturb the input at a given position, can the output change at all? Posing this question one input
at a time yields a vector of Boolean partial ``derivatives'', and for a program with several outputs, a
\emph{Boolean Jacobian}: a matrix whose entry at $(j, i)$ records whether output position $j$ may depend on
input position $i$. Consider multiplication: at the point $(x, y)$, the usual partial derivative $\partial(x
\cdot y)/\partial x$ is $y$, and its Boolean counterpart records that the output can vary with $x$ only when
$y \neq 0$. Like their numerical counterparts, Boolean Jacobians compose by matrix multiplication, with
conjunction and disjunction playing the role of multiplication and addition. This paper presents an approach
to dynamic dependency analysis where such Jacobians are evaluated alongside the program, forwards or
backwards, in the manner of forward- and reverse-mode AD, with the two directions related by matrix
transposition.

Existing approaches to program slicing are often tied to particular programming languages or implementations.
In this paper we develop this analogy into a general categorical approach to data provenance, in
which the derivative of a program is a linear map between spaces of approximations of its input and output,
with scalars drawn from a commutative semiring of dependency information. The choice of semiring determines
the analysis: the reals recover AD itself, while over distributive lattices the derivative and its transpose
form a \emph{conjugate pair}, propagating dependency information forwards and backwards; when the lattices are
Boolean algebras, the same maps can instead be presented as an adjunction, recovering the \GPS of Perera and
collaborators. Our categorical approach should provide a suitable setting for enabling ``automatic''
data provenance for a variety of programming languages, and is easily configured to use alternative
approximation strategies, including quantitative forms of slicing.

\subsection{Dependencies as Derivatives}
\label{sec:introduction:galois-slicing}

Consider a single run of a program on a particular input. Each part of the output is computed from
certain parts of the input, and recording these facts gives the \emph{dependency relation} of the run,
relating each output position to the input positions it depends on. This section works through a small example
showing how dependency relations behave as derivatives: they have a tabular representation as Boolean
matrices, compose using matrix multiplication, and can be evaluated forwards or backwards.

\begin{example}
  \label{ex:introduction-example}
  The following program is written in Haskell syntax \cite{haskell}, using a list comprehension to filter a list of pairs of labels and numbers to those numbers with a given label, and then computing the sum of the numbers:
  \begin{displaymath}
    \begin{array}{l}
      \mathrm{query} :: \mathrm{Label} \to [(\mathrm{Label}, \mathrm{Int})] \to \mathrm{Int} \\
      \mathrm{query}\,l\,\mathit{db} = \mathrm{sum}\,[ n \mid (l',n) \leftarrow \mathit{db}, l \equiv l' ]
    \end{array}
  \end{displaymath}
  With $\mathit{db} = [(\mathsf{a}, 0), (\mathsf{b}, 1), (\mathsf{a}, 1)]$, we will have $\mathit{query}\,\mathsf{a}\,\mathit{db}$ and $\mathit{query}\,\mathsf{b}\,\mathit{db}$ both evaluating to $1$.

  Suppose we are interested in how the output depends on the numerical parts of the input, for the
  query parameters $l = \mathsf{a}$ and $l = \mathsf{b}$. The three numbers in $\mathit{db}$
  occupy three \emph{positions}, and the output a single position. In the run with $l = \mathsf{a}$, the
  output is computed from the numbers at the first and third positions; in the run with $l = \mathsf{b}$,
  from the number at the second. We depict each input position pointing to the output positions that depend
  on it:
      \begin{center}
    \begin{tikzpicture}[every node/.style={font=\small}]
      % query a
      \node (la) at (0,1.4) {$\mathit{query}\;\mathsf{a}$};
      \node (a1) at (0,0.8) {$(\mathsf{a},0)$};
      \node (a2) at (0,0) {$(\mathsf{b},1)$};
      \node (a3) at (0,-0.8) {$(\mathsf{a},1)$};
      \node (ao) at (2.4,0) {$1$};
      \draw[->] (a1) -- (ao);
      \draw[->] (a3) -- (ao);
      % query b
      \node (lb) at (6,1.4) {$\mathit{query}\;\mathsf{b}$};
      \node (b1) at (6,0.8) {$(\mathsf{a},0)$};
      \node (b2) at (6,0) {$(\mathsf{b},1)$};
      \node (b3) at (6,-0.8) {$(\mathsf{a},1)$};
      \node (bo) at (8.4,0) {$1$};
      \draw[->] (b2) -- (bo);
    \end{tikzpicture}
  \end{center}
    Tabulating each relation, with a row for each output position and a column for each input position,
  gives the \emph{Jacobian} of the run:
    \begin{displaymath}
    \partial(\mathit{query}\,\mathsf{a}\,\mathit{db}) =
    \begin{pmatrix} 1 & 0 & 1 \end{pmatrix}
    \qquad
    \partial(\mathit{query}\,\mathsf{b}\,\mathit{db}) =
    \begin{pmatrix} 0 & 1 & 0 \end{pmatrix}
  \end{displaymath}
    As in differential calculus, the Jacobian is taken at a point: the two runs of the same program have
  different Jacobians. And like their numerical counterparts, these matrices compose: $\mathit{query}\,l$ is a
  sum applied to a selection, and the Jacobian of the run is the product of the Jacobians of those two stages,
  a chain rule for dependencies.

  In setting up the positions, we chose to make only the numbers perturbable, keeping the labels and
  the structure of the list itself fixed. Other choices are also useful, and one of the aims of this work is
  to clarify how to choose a dependency structure appropriate for different tasks by selecting an appropriate
  semiring of dependency information. We elaborate on this further in
  \secref{semantic-gps}.

  The Jacobian acting directly evaluates dependencies \emph{forwards}: given a selection of input
  positions, it returns the output positions that depend on them. A selection of positions can be displayed as
  an \emph{approximation} of the underlying value, with the numbers at unselected positions replaced by
  $\bot$. Writing $\partial(\mathit{query}\,l\,\mathit{db})_f$ for this forward action, we have
  $\partial(\mathit{query}\,\mathsf{a}\,\mathit{db})_f([(\mathsf{a},0),(\mathsf{b},\bot),(\mathsf{a},\bot)]) =
  1$, since the output depends on the entry $(\mathsf{a},0)$, whereas
  $\partial(\mathit{query}\,\mathsf{b}\,\mathit{db})_f$ applied to the same selection returns $\bot$: that run
  consulted no $\mathsf{a}$-labelled entries.

  The transpose of the Jacobian evaluates dependencies \emph{backwards}: given a selection of
  output positions, it returns the input positions on which they depend. Writing
  $\partial(\mathit{query}\,l\,\mathit{db})_r$ for this reverse action, the fully approximated output $\bot$
  requires none of the input:
  \begin{displaymath}
    \partial (\mathit{query}\,l\,\mathit{db})_r(\bot) = [(\mathsf{a},\bot), (\mathsf{b}, \bot), (\mathsf{a}, \bot)]
  \end{displaymath}
  for both $l = \mathsf{a}$ and $l = \mathsf{b}$. If instead we retain the output ``$1$''
  unapproximated, then the two
  query runs' backwards actions return different results:
  \begin{displaymath}
    \begin{array}{l}
      \partial (\mathit{query}\,\mathsf{a}\,\mathit{db})_r(1) = [(\mathsf{a},0), (\mathsf{b},\bot), (\mathsf{a},1)] \\
      \partial (\mathit{query}\,\mathsf{b}\,\mathit{db})_r(1) = [(\mathsf{a},\bot), (\mathsf{b},1), (\mathsf{a},\bot)]
    \end{array}
  \end{displaymath}
  Pieces of the input that were {\em not} used are replaced by $\bot$. As we expect, the run of the query with label $\mathsf{a}$ depends on the entries in the database labelled with $\mathsf{a}$, and likewise for the run with label $\mathsf{b}$.

  In a simple query like this, it is easy to work out the dependency relationship between the input and output. However, the benefit of language-based approaches is that they are {\em automatic} for all programs, no matter how complex the relationship between input and output. Moreover, by changing what we mean by ``approximation'' we can compute a range of different information about a program.
\end{example}

\subsection{Galois Slicing and Automatic Differentiation}

The backward maps of the previous section compute a form of \emph{data provenance}: for each part of
the output, the part of the input needed to explain it. Dependency information of this kind is the common
currency of a range of existing techniques, from program slicing~\cite{weiser81} to provenance for database
queries~\cite{buneman95,cheney07}, including semiring provenance~\cite{green07}. Closest
to our setting is the \GPS of Perera and
collaborators~\cite{perera12a,perera16d,ricciotti17}, which organises the forward and backward maps of a run
into a Galois connection between lattices of approximations.

Many forms of dynamic provenance analysis, \GPS included, work by recording a trace of each
execution and computing the backwards analysis by re-running over the trace. A denotational account would
instead bake the analysis into the semantics of the program itself, rather than provide it as a separately
defined ``backwards evaluation'' operation. 
The differential reading of dependency analysis, developed in \secref{approx-as-tangents}, points to
a way to obtain such an account, by analogy with \emph{automatic
differentiation}~\cite{siskind08,elliott18,vakar22}; let us make the correspondence explicit.

\begin{itemize}[leftmargin=\enummargin]
\item For dependency analysis, every value has an associated join-semilattice of \emph{approximations}, with least element $\bot$ (the empty selection). For differentiable programs, every point has an associated vector space of {\em tangents}.
\item For dependency analysis, every program has an associated forward approximation map that takes approximations of the input to the approximations of the output that depend on them. This map {\em preserves joins and $\bot$}. For differentiable programs, every program has a forward derivative that takes tangents of the input to tangents of the output. The forward derivative map is {\em linear}, so it preserves addition of tangents and the zero tangent.
\item For dependency analysis, every program has an associated backward approximation map that takes approximations of the output back to least approximations of the input that they depend on. This map also {\em preserves joins and $\bot$}. For differentiable programs, every program has a reverse derivative that takes \emph{cotangents} of the output (linear maps from tangents to scalars) to cotangents of the input. This map is again {\em linear}.
\item In both settings, the forward and backward maps are related by being each other's transpose.
The transpose makes sense because the spaces are \emph{self-dual}: tangent vectors pair with tangent vectors
via the dot product, identifying cotangents with tangents, and selections of positions pair by the same dot
product formula, with joins and meets in place of sums and products.
\end{itemize}

Given this close connection between dependency analysis and differentiable programming, we can take structures intended for modelling automatic differentiation, such as V{\'a}k{\'a}r's CHAD framework and use them to model dependency analysis. This will enable us to generalise and expand the scope of automatic differentiation to act as a foundation for data provenance in a wider range of computational settings.

\subsection{Outline and Contributions}

%Our primary contribution in this article is to elucidate \GPS by relating it to differentiable programming and automatic differentiation. In doing so, we aim to expand the applicability of program slicing, and to potentially transfer efficient automatic differentiation implementation techniques from automatic differentiation to program slicing and provenance tracking.

Our main contribution is to show that data provenance can be understood as differentiation over a
commutative semiring of dependency information. In \secref{approx-as-tangents} we develop the first-order
picture, for programs over tuples of atomic data: derivatives as matrices over a semiring $S$, with the
transpose relating forward and backward dependency analysis, and with successively stronger conditions on $S$ yielding
join-preserving maps, conjugate pairs and, for Boolean algebras, the Galois connections of \GPS; ordinary AD
occupies the case $S = \mathbb{R}$. We also relate this picture to \citet{berry79}'s \emph{stable domain
theory}, where stable functions provide a kind of semantic provenance analysis.

The following table summarises the correspondence developed over the course of the paper:
\vspace{1.5mm}
\begin{center}
  \small
  \begin{tabular}{lll}
    \toprule
    & \textbf{Differentiable programming} & \textbf{Dependency analysis} \\
    \midrule
    scalars & real numbers & commutative semiring $S$ \\
    space at a point & vector space of tangents & semimodule of approximations over $S$ \\
    identification with the dual & finite dimensionality & chosen self-duality \\
    derivative at a point & Jacobian matrix & $S$-weighted relation as matrix over $S$ \\
    chain rule & matrix multiplication & matrix multiplication over $S$ \\
    forward mode & push tangents forward & propagate dependencies forward \\
    reverse mode & pull cotangents back & propagate dependencies backwards \\
    forward/backward relation & transpose & transpose via self-dualities \\
    \bottomrule
  \end{tabular}
\end{center}
\vspace{1.5mm}

In \secref{models-of-total-gps} we extend to sum types using the category of families construction and,
following V{\'a}k{\'a}r \etal{}'s CHAD framework~\cite{vakar22,nunes2023}, build models of a higher-order
language in which every program is interpreted together with its family of derivatives. We apply this to a
concrete higher-order language in \secref{language} and demonstrate the model on a number of examples, showing
how parameterising on $S$ controls the dependency structure associated with data, something that was
``hard-coded'' in previous presentations of \GPS. We prove two correctness properties in
\secref{definability}: programs of first-order type denote matrices in the sense of
\secref{approx-as-tangents}, and the higher-order interpretation agrees with the standard set-theoretic
semantics. \secref{related-work} and \secref{conclusion} discuss additional related and future work.

% \begin{enumerate}[leftmargin=\enummargin]
% \item We explain how the CHAD framework of Vákár et al. can be adapted to give a general categorical framework for \GPS.
% \item Using our adapted CHAD framework, we can explain how type structure can be used to control the approximation lattices associated with data points, something that was ``hard coded'' in previous presentations of \GPS.
% \item With the benefit of our abstract setting, we can relate \GPS to other parts of denotational semantics. In particular, we show that there is a close connection with Stable Domain Theory, a proposed solution to capturing sequentiality by recording more intensional information about programs' sensitivity to approximations.
% \end{enumerate}

We have formalised our major results in Agda, resulting in an executable implementation built directly from the categorical constructions that we have used to compute the examples in \secref{language:examples}. Please consult the file \texttt{everything.agda} in the supplementary material.

% \begin{enumerate}
% \item We explain how the CHAD framework of Vákár et al. can be adapted to give a general categorical framework for \GPS.
% \item Using our adapted CHAD framework, we can explain how type structure can be used to control the approximation lattices associated with data points, something that was ``hard coded'' in previous presentations of \GPS.
% \item With the benefit of our abstract setting, we can relate \GPS to other parts of denotational semantics. In particular, we show that there is a close connection with Stable Domain Theory, a proposed solution to capturing sequentiality by recording more intensional information about programs' sensitivity to approximations.

% We have \anonymise{\href{https://github.com/bobatkey/approx-diff}{formalised this work in Agda}}{formalised this work in Agda}. Not only does this mean that the constructions and proofs have been checked, but also the construction is executable and can be run on small examples. We are also planning to use Agda's JavaScript backend to produce a demo.

\section{Approximations as Tangents}
\label{sec:approx-as-tangents}

We motivate our approach by showing how to combine ideas from differential geometry and linear algebra to
reconstruct the dependency analyses of \secref{introduction} in a denotational setting.

\subsection{Manifolds, Smooth Functions, and Automatic Differentiation}
\label{sec:approx-as-tangents:autodiff}

% For what follows we do not need to know in detail what a manifold is, or exactly how smooth functions are defined. We only state some definitions in order to fix terminology and to justify our claim of a connection to manifolds and smooth functions. \bob{Need a good reference here.}

The general study of differentiable functions takes place on \emph{manifolds}, topological spaces that
``locally'' behave like an open subset of the Euclidean space $\RR^n$. The spaces $\RR^n$ themselves are
manifolds, but so are ``non-flat'' examples such as $n$-spheres and yet more exotic spaces. Every point $x$ in
a manifold $M$ has an associated \emph{tangent vector space} $\tangents_x(M)$ consisting of linear
approximations of curves on the manifold passing through $x$. Each point also has a \emph{cotangent vector
space} $\cotangents_x(M) = \tangents_x(M) \linearto \RR$. The tangent and cotangent spaces are finite
dimensional, so in the presence of a chosen basis they are canonically isomorphic. In the case when the
manifold is $\RR^n$, then every tangent space is isomorphic to $\RR^n$ as well.

Smooth functions $f$ between manifolds $M$ and $N$ are functions on their points that are locally
differentiable on $\RR^n$. Manifolds and smooth functions form a category $\Man$. Each smooth function induces
maps of the (co)tangent spaces:
\begin{itemize}[leftmargin=\enummargin]
\item The \emph{forward derivative} (tangent map, pushforward) $\pushf{f}_x$ is a linear map $\tangents_x(M) \linearto \tangents_{f(x)}(N)$. In the Euclidean case when $M = \RR^m$ and $N = \RR^n$, the tangent map can be represented by the Jacobian matrix of partial derivatives of $f$ at $x$.
\item The \emph{backward derivative} (cotangent map, pullback) $\pullf{f}_x$ is a linear map $\cotangents_{f(x)}(N) \linearto \cotangents_x(M)$. In the Euclidean case, the backward derivative is represented by the transpose of the Jacobian of $f$ at $x$.
\end{itemize}

\begin{remark}[Chain Rule]
  \label{rem:chain-rule}
  A useful property of derivative maps is that they compose according
  to the chain rule. Suppose that $f : M \to N$ and $g : N \to K$ are
  smooth functions. Then for any $x \in M$, we have:
  \begin{itemize}
  \item $\pushf{(g \circ f)}_x = \pushf{g}_{f(x)} \circ \pushf{f}_x : \tangents_x(M) \linearto \tangents_{g(f(x))}(K)$
  \item $\pullf{(g \circ f)}_x = \pullf{f}_x \circ \pullf{g}_{f(x)} : \cotangents_{g(f(x))}(K) \linearto \cotangents_x(M)$
  \end{itemize}
  The chain rule has the practical effect that we can compute
  derivative maps of $f$ and $g$ independently and compose them,
  instead of the potentially more difficult task of computing the
  derivative maps of $g \circ f$. As we shall see below, stable maps
  also obey a chain rule, and this forms the basis of the general
  categorical approach to differentiability that we describe in
  \secref{models-of-total-gps}.
\end{remark}

Computing the forward and backward derivatives of smooth functions $f$ has many applications of practical
interest. For example, computation of the reverse derivative is of central interest in machine learning by
gradient descent, the main technique used to train deep neural networks~\cite{rumelhart88,goodfellow16}.

Derivatives can be computed numerically by computing $f$ on small perturbations of its input, or symbolically
by examining a closed-form representation of $f$. However, a more common and practical technique is to use
\emph{automatic differentiation}, where a program computing $f$ is instrumented to produce (a representation
of) the forward and/or backward derivative as a side-effect of producing the output~\cite{linnainmaa76}. This
has led to the area of differentiable programming, where programming languages and their implementations are
specifically designed to admit efficient automatic differentation
algorithms~\cite{jax2018github,abadi16,elliott17,sigal24}.

\subsection{\added{Potential Change as Tangents}}

\Added{Differential Calculus tracks how functions adapt to changes in
  their inputs. The tangent space at a point is the vector space of
  all potential changes applicable at that point. The forward and
  reverse derivatives transport changes forward and back. We track
  data provenance using a simplification of this idea where tangents
  are taken from the Boolean semiring $\{0,1\}$ where $0$ represents
  ``no change'' and $1$ represents ``possible change''. Potential
  change to an $n$-tuple of data is represented as an $n$-vector of
  $\{0,1\}$ values. These vectors form a semimodule over the semiring
  structure on $\{0,1\}$ with maximum (join) as addition and minimum
  (meet) as multiplication. Thus ``no change'' is represented as the
  $0$ vector $(0 \dots 0)$.}

\Added{With this definition of ``tangent'', the derivative of a
  function
  $f : A_1 \times \cdots \times A_m \to B_1 \times \cdots \times B_n$
  ought to map tangent vectors in $\{0,1\}^m$ to tangent vectors in
  $\{0,1\}^n$, describing how potential changes in the input are
  reflected in the output. Differentiable functions on manifolds have
  an intrinsic notion of derivative, but for plain functions $f$ we
  require that they are paired with a tangent map that describes the
  intensional information about how they depend on their
  arguments. For $f$, the tangent map $\partial f$ has type
  $A_1 \times \cdots \times A_n \to \{0,1\}^{m\times n}$, mapping
  points to matrices representing linear maps: the ``Jacobian'' for
  this function. We are not free to choose an arbitrary linear map for
  $\partial f(a_1, \dots, a_m)$, but must choose one that satisfies
  the following \emph{dependency correctness} property:}

\begin{AddedBlock}
  \begin{definition}[Dependency Correctness] A tangent map
    $\partial f : \overrightarrow{A_i} \to \{0,1\}^{m \times n}$ is
    \emph{dependency correct} for a function
    $f : \overrightarrow{A_i} \to \overrightarrow{B_j}$ if for all
    $a, a' \in \overrightarrow{A_i}$ and $0 \leq j < m$, if $a$ and
    $a'$ are equal on all components $i$ such that
    $\partial f(a)_{i,j} = 1$, then $f(a)_j = f(a')_j$.
  \end{definition}
\end{AddedBlock}

\Added{Reading the definition backwards, this says that $f$ is
  constant at component $j$ on inputs that differ only at components
  $i$ where the tangent map matrix says $0$. Thus we read the entry at
  $(i,j)$ of the tangent matrix as capturing the dependency of output
  $j$ on input $i$. We illustrate this with an example:}

\begin{AddedBlock}
  \begin{example}
    Let $\pi_1 : A_1 \times A_2 \to A_1$ be first projection. At all
    points $(a_1, a_2)$, the tangent matrix is $(1\;0)$ indicating
    that the output depends on the first argument but not the
    second. This is justified by observing how it acts on tangent
    vectors. Representing no change in the first argument and
    potential change in the second as
    $\left(\begin{array}{c} 0 \\ 1 \end{array}\right)$, applying the
    tangent map gives $0$ indicating no change in the output. Swapping
    the potential change to
    $\left(\begin{array}{c} 1 \\ 0 \end{array}\right)$ yields the
    output tangent $1$ indicating that changing the first argument
    changes the output. By linearity, all other potential changes are
    derivable from these. In particular, for all tangent maps, no
    change in the input ($\overrightarrow{0}$) is always mapped to no
    change in the output.
  \end{example}
\end{AddedBlock}

\Added{Note that dependency correct tangent maps for a function are
  not unique. We can always overapproximate the dependency of a
  function on its input. In the case of a binary function like first
  projection, $\partial\pi_1(a_1,a_2) = (1\;1)$ is also permissble.  }

\begin{AddedBlock}
  \begin{example}
    Multiplication $(*) : \mathbb{N} \times \mathbb{N} \to \mathbb{N}$
    is an example of a function that has tangent maps that depend on
    the input. Dependency correct tangent maps can be defined as:
    \begin{displaymath}
      \begin{array}{lcl}
        \partial(*)(0,0)&=&(1\;1) \\
        \partial(*)(0,n)&=&(1\;0) \\
        \partial(*)(m,0)&=&(0\;1) \\
        \partial(*)(m,n)&=&(1\;1)
      \end{array}
    \end{displaymath}
    These maps indicate that the dependency of multiplication on an
    input depends on whether or not other other input is $0$. If the
    first argument is zero, then changing the second will have no
    effect and vice versa. If both arguments are $0$, then a change in
    both will change the output. However, since tangent maps are
    componentwise we have no way of saying change in both, so we
    overapproximate by saying that a change in either \emph{may}
    change the output. Again, this choice is not unique. We could also
    use $(1\;0)$ or $(0\;1)$ and still be dependency correct,
    indicating a choice to be left- or right-biased.
  \end{example}
\end{AddedBlock}

\Added{Functions with dependency correct tangent maps have an identity
  and compose according to the chain rule, so they form a category.}

\newcommand{\DepCat}{\mathbf{DepCat}}

\begin{AddedBlock}
  \begin{proposition}
    The identity tangent map is dependency correct for the identity
    function. Composable functions $f, g$ with dependency correct
    tangent maps $\partial f$ and $\partial g$ are closed under
    composition, with the tangent map being
    $\partial(f \circ g)(x) = \partial f (g(x)) \circ \partial g(x)$.

    The category $\DepCat$ of functions with dependency information
    has objects that are tuples of sets $\overrightarrow{A_i}$ and
    morphisms that are functions with dependency correct tangent
    maps. This category has finite products.
  \end{proposition}

  FIXME: finite products is interesting, because it determines the
  dependency information???

  The fact that dependency correct functions form a category allows us
  to model a simple programming language with finite product
  types. Primitive operations can be interpreted a functions that
  overapproximate their dependency information:

  \begin{proposition}
    There is a functor $\Set \to \DepCat$ that preserves finite
    products. FIXME: check!
  \end{proposition}

  Morphisms in $\DepCat$ constructed from finite product structure and
  embedded functions will only have trivial tangent maps arising from
  the use of projections. A generator of non-trivial dependency
  information is the ability to use parts of the input depending on
  the input itself, via a condition (if-then-else) map:

  \begin{proposition}
    Let $\mathbb{B}$ be the object of $\DepCat$ consisting of the set
    of booleans $\{\mathsf{true}, \mathsf{false}\}$. There are
    dependency correct maps
    $\mathrm{true}, \mathrm{false} : 1 \to \mathbb{B}$ and a natural
    family of maps
    $\mathrm{cond}_X : \mathbb{B} \times (X \times X) \to X$ such that $\mathrm{cond}_X \circ \langle \mathrm{true} \circ !, \mathrm{id} \rangle = \pi_1$ and $\mathrm{cond}_X \circ \langle \mathrm{false} \circ !, \mathrm{id} \rangle = \pi_2$. FIXME: and extensionality.
  \end{proposition}

  Thus $\DepCat$ is suitable for modelling a small language with
  booleans, conditionals, and a collection of primitive types with
  operations. This setting does not cover our example from the
  introduction, which used lists and higher-order functions. In
  Section \ref{sec:models-of-total-gps} we embed $\DepCat$ into
  categories that account for sum types, some algebraic types and
  higher-order functions.
\end{AddedBlock}

\subsection{Derivatives over a Commutative Semiring}
\label{sec:approx-as-tangents:semiring}

In the Euclidean case, the differential picture above is plain matrix algebra: with a chosen basis, tangent
and cotangent spaces are both $\RR^n$, the forward derivative at a point is the Jacobian matrix, and the
backward derivative is its transpose. Once the picture is presented in coordinates like this, the real numbers
play no further role beyond supplying the scalars; since dependency information is typically indexed by
positions in a value, which then serve as a coordinate basis, our first move is to keep the matrix algebra and
change the scalars. Thus throughout this section, we fix a commutative semiring $S = (S, +, 0, \cdot, 1)$,
whose elements we think of as atoms of dependency information, with $\cdot$ combining information along a path
from an input to an output and $+$ combining information across parallel paths.

Our second move is one of abstraction: rather than axiomatise the matrices of \secref{introduction} directly,
we identify the minimal structure required to capture what made that story work: a space of dependency
information for each value, linear maps between spaces for the forward analysis, and an identification of each
space with its dual, the extra structure needed to support an analogue of matrix transposition, and hence a
backward analysis. The role of vector spaces, which are defined over fields, is played by \emph{semimodules}:
commutative monoids $(M, +, 0)$ equipped with an action $a \cdot x$ of the scalars, satisfying the usual laws.
Linear maps preserve $+$, $0$ and the action, and form a category $\SemiMod(S)$; pointwise addition of linear
maps makes each hom-set a commutative monoid, with composition bilinear. Every semimodule has a \emph{dual}
$M^* = M \linearto S$ of linear maps into the scalars (categorically, the internal hom into the scalars
object, which is the unit of the tensor product of semimodules), and every linear map $f : M \linearto N$ has
a \emph{transpose} $\transpose{f} = (- \comp f) : N^* \linearto M^*$, contravariantly.

The transpose is almost the backward analysis we are after: it runs in the opposite direction to $f$, but
between the duals $N^*$ and $M^*$, the cotangent spaces, rather than between the tangent spaces $M$ and $N$
themselves. What we want is a map $N \linearto M$, running dependency information back through the same spaces
the forward map runs it through, as the transposed matrix did in \secref{introduction}. For that, each space
must be identified with its dual. Whereas a finite-dimensional vector space always admits such an
identification (choose a basis), a semimodule need not; we therefore make the identification part of the data.
\begin{definition}[Self-dual semimodule]
  \label{def:self-dual-semimodule}
  A \emph{self-dual} semimodule is a pair $\langle M, i \rangle$ of a semimodule and a chosen isomorphism
  $i : M \cong M^*$. The \emph{pairing} induced by $i$ is $\langle x, y \rangle = i(x)(y)$.
\end{definition}
The pairing is an $S$-valued measure of the extent to which $x$ and $y$ overlap; its prototype is the
dot product of vectors (the pairing induced by the self-duality of $S^n$, as we will see below). Extending the
analogy, we call $x$ and $y$ \emph{orthogonal} when $\langle x, y \rangle = 0$.
\begin{definition}[Conjugate]
  \label{def:conjugate}
  Between self-dual semimodules $\langle M, i \rangle$ and $\langle N, j \rangle$, the \emph{conjugate} of
  a linear map $f : M \linearto N$ is $\conj{f} = i^{-1} \comp \transpose{f} \comp j : N \linearto M$, the
  unique linear map satisfying
  \begin{displaymath}
    \langle f(x), y \rangle = \langle x, \conj{f}(y) \rangle
  \end{displaymath}
\end{definition}
The conjugate is our backward analysis, and its defining equation makes precise the relationship with
the forward one alluded to in \secref{introduction}. Self-dual semimodules and linear maps form a category
$\SDSemiMod(S)$.

The matrices of \secref{introduction} now reappear as concrete instances. The \emph{free} semimodule $S^n$ on
$n$ generators has as elements the $n$-length vectors of scalars, with standard basis vectors $e_i$ and every
$v \in S^n$ decomposing as $v = \sum_i v_i \cdot e_i$. Currying the \emph{dot product} $\langle u, v \rangle =
\sum_i u_i \cdot v_i$ gives an isomorphism $S^n \cong (S^n)^*$, since pairing with the basis vectors extracts
components, $\langle e_i, v \rangle = v_i$; so free semimodules are self-dual. By basis decomposition, a
linear map $S^m \linearto S^n$ is exactly an $n \times m$ matrix over $S$, with composition as matrix
multiplication and the conjugate as the familiar matrix transpose, $(\conj{M})_{ji} = M_{ij}$.

We write $\Mat(S)$ for the skeletal category whose objects are the natural numbers and whose morphisms $m \to
n$ are the $n \times m$ matrices. The Boolean Jacobians from \secref{introduction} live here: a dependency
relation between $m$ input positions and $n$ output positions is exactly a morphism $m \to n$ in $\Mat(\Two)$,
where $\Two$ is the Boolean semiring with $+ = \vee$ and $\cdot = \wedge$, and the chain rule for dependencies
is composition in $\Mat(\Two)$, or equivalently \emph{relational} composition under the equivalence of
categories $\Mat(\Two) \simeq \FinRel$. The dot product can be interpreted as ``overlap'', becauses $\langle
u, v \rangle = \top$ exactly when the selections $u$ and $v$ share a position, i.e.\ are not disjoint.

The trivial semimodule is a zero object in $\SDSemiMod(S)$, and self-dual semimodules are closed under the
\emph{biproduct} $M \biprod N$, an object which is both a product and a coproduct, with the injections
conjugate to the projections, $\biinj_i = \conj{\biproj_i}$. The coincidence of products and coproducts is
characteristic of such linear settings, and means $\SDSemiMod(S)$ on its own cannot interpret programs over
sum types; we return to this in \secref{models-of-total-gps}.

\subsection{Dependency Structures over Semirings}
\label{sec:approx-as-tangents:richer-semirings}

The framework so far only asks for commutativity of the scalars, which makes the dual $M^*$ an $S$-semimodule
again, so that the transpose remains a map between $S$-semimodules. Any commutative semiring $S$ thus yields a
category $\SDSemiMod(S)$ with forward derivatives and their conjugates, related through the pairing. When $S$
is the reals, the conjugate is the familiar adjoint of linear algebra. In a dependency-tracking setting,
however, the relationships of interest are typically granular: we ask whether, or to what degree, an output
depends on an input, and want to compare one answer with another. This calls for an order on dependency
information, which additional axioms on $S$ provide.

\subsubsection{Conjugate pairs and conjugacy analyses}

Suppose first that $1$ is an additive top in $S$ ($1 + a = 1$). Addition is then
idempotent, meaning $a + a = a$. Idempotent addition induces an order: $S$ becomes a join-semilattice,
with $a \leq b$ iff $a + b = b$, joins given by $+$ and least element $0$. This structure lifts pointwise to every semimodule,
and every linear map automatically preserves joins and $0$. Dependency
information is now ordered by informativeness: $0$ means ``no dependency'', and the forward and backward
analyses are monotone, taking more informative inputs to more informative outputs and vice-versa.

If multiplication is also idempotent ($a \cdot a = a$), then $S$ is a bounded distributive lattice,
with multiplication as the meet and distributivity given by the semiring. This opens the way to reading the
two analyses in terms of \emph{disjointness}: where meets exist, define $x \disj y$ when $x \meet y = 0$, and
if disjointness agrees with the orthogonality given by the pairing, the defining equation of the conjugate
(\defref{conjugate}) says that $f$ and $\conj{f}$ exchange disjointness: $f(x) \disj y$ iff $x \disj
\conj{f}(y)$. Let us therefore consider the case where each semimodule has meets, preserved by the scalar action and
agreeing with orthogonality in this sense; meets, unlike joins, do not lift from $S$, so this is additional
structure. Maps that exchange disjointness in this way have a standard name.
\begin{definition}[Conjugate pair]
  \label{def:conjugate-pair}
  Join-preserving maps $f : L \to M$ and $g : M \to L$ between bounded lattices form a \emph{conjugate
  pair} when $f(x) \disj y$ iff $x \disj g(y)$ for all $x \in L$ and $y \in M$.
\end{definition}
The notion originates in relation algebra~\cite{jonsson51}, where the conjugate of composition with a
relation is composition with its converse. The scalars $S$, viewed as a semimodule over itself, qualify
automatically for this reading, and biproducts (including the empty biproduct $S^0$) inherit the structure
componentwise, giving the linear maps between free semimodules $S^n$, along with their conjugates, as the canonical
instances of conjugate pairs.

The relevance to us here is that conjugate pairs also underpin the \emph{cognacy} (common ancestry)
analyses of work on linked visualisations~\cite{perera22,bond25}. These dependency analyses work by composing
a join-preserving dependency analysis with its conjugate, yielding a map sending a selection of outputs to its
\emph{related outputs}, those sharing a dependency on some part of the input; selections in one chart can then
be linked to ``cognate'' selections in another.

\subsubsection{Galois slicing}

Suppose in addition that every scalar has a complement: an operation $\neg$ with $a \meet \neg a =
0$ and $a \join \neg a = 1$, making $S$ a Boolean algebra. Negation supports an alternative presentation of
conjugates as adjoints. Composing a map with complements on either side preserves its direction but exchanges
join-preservation for meet-preservation; the \emph{De Morgan dual} $\neg \comp \conj{f} \comp \neg$ of
the conjugate is thus meet-preserving, and is precisely the upper adjoint of $f$ (upper adjoints
preserve meets, dually to lower adjoints and joins). Forward and backward analysis therefore form a
\emph{Galois connection} between the lattices of approximations.
\begin{definition}[Galois connection]
  \label{def:galois-connection}
  Suppose $X$ and $Y$ are posets. A \emph{Galois connection} $f \adj g: X \to Y$ is a pair of monotone
  functions $f: Y \to X$ and $g: X \to Y$ satisfying $y \leq g(x) \iff f(y) \leq x$ for any $x \in X$ and $y
  \in Y$. Since a Galois connection is also an adjunction, we refer to $f$ as the left adjoint and $g$ as the
  right adjoint.
\end{definition}
This is precisely the setting of \GPS~\cite{perera12a,perera16d,ricciotti17}, recovered here as the
Boolean instance of the general semimodule picture. Galois connections as such presuppose none of the
algebraic structure of this section, being definable between arbitrary posets; De Morgan duality is simply how
the adjoint presentation arises here, where the conjugate is the primary notion. One orientation to keep
straight: the conjugate as presented here, $\conj{f}$, runs backwards, whereas the meet-preserving map of \GPS
runs forwards. The two presentations agree because $\conj{-}$ is involutive: instantiating the construction at
$\conj{f}$ in place of $f$ gives exactly the connection of \GPS, with meet-preserving forward map $\neg \comp
f \comp \neg$.

The join-preserving and meet-preserving maps answer different questions. In the forwards direction, the
join-preserving map $f$ has a \emph{necessity}-flavoured interpretation: which outputs may depend on a given
part of the input. Its meet-preserving counterpart $\neg \comp f \comp \neg$ reads as \emph{sufficiency}:
which outputs a given input selection determines outright. Notably, in the Galois slicing work the
meet-preserving forward map is only used to establish the existence of the join-preserving backward map; it
seems to lack any direct dependency analysis applications of its own.

\subsubsection{Perturbation bounds}
\label{sec:approx-as-tangents:perturbation-bounds}

Dependency analysis over the Booleans records whether an output can change when an input does;
\exref{introduction-example} raised the more fine-grained question of how much it
can change. Numerical analysis studies this question as \emph{propagation of error}: given bounds on how far
each input may be perturbed, derive a bound on the output~\cite{high:ASNA2}. Propagation of error is again a
choice of semiring, this time the \emph{tropical semiring} $(\mathbbm{Q}_{\geq 0} \cup \{\infty\}, \min,
\infty, +, 0)$~\cite{simon88}. A scalar is a non-negative rational (or $\infty$) bounding the extent to which
a value may change from its value in the run, with $\infty$ meaning unconstrained. Semiring addition is
$\min$, so that combining bounds for parallel paths keeps the tighter bound, and multiplication is numerical
$+$, accumulating bounds along a path. Addition is idempotent, so the analyses are monotone as in
\secref{approx-as-tangents:richer-semirings}: smaller bounds are more informative, and the semiring zero
$\infty$ means ``no information''.

A Jacobian entry $c$ now acts as a translation: if the input at that position changes by at most
$\delta$, with the others held fixed, the output changes by at most $\delta + c$. Addition's entries are $0$: a
bound on either argument passes through undisturbed. Multiplication, however, scales perturbations: if $x$
changes by $\delta$ then $x \cdot y$ changes by $|y|\,\delta$, and no translation bounds a scaling; its
entries are $\infty$. This is sound, but reports nothing: propagation of absolute error is blind to
multiplication.

Measuring perturbations multiplicatively instead keeps $\min$ as the join of bounds but accumulates
along paths by numerical product: $(\mathbbm{Q}_{\geq 0}
\cup \{\infty\}, \min, \infty, \times, 1)$, with scalars now read as factors. This is the
multiplicative counterpart of the previous presentation, related to it by logarithms over the reals, and the
two are the classical dichotomy between absolute and relative error~\cite[\S\S 1.6--1.7]{high:ASNA2}.
Multiplication now has unit entries, since relative perturbations pass through a product unchanged, while an
addition $x + y$ acquires the entry $|x| / |x + y|$: the factor by which it amplifies relative error, its
(relative) \emph{condition number}.

In the refund example the final sum adds the products $6$ and $-6$, so the amplification is $|6| / |6 +
(-6)|$: an infinite condition number, known in numerical analysis as catastrophic cancellation. For
non-negative data the factors are at most $1$, so sums never amplify; the refund's negative entry makes the
infinite factor possible. Where the Boolean analysis could only report a dependency that was not really
there, the relative analysis says how badly the bound degrades. The difference between the two is a choice of
semiring, and making that choice is part of specifying what the analysis is to observe.

\subsubsection{Qualitative derivatives}

The Boolean analysis reports a dependency at a cancellation because it forgets signs. The \emph{sign
semiring} $\{+, 0, -, ?\}$, the rule-of-signs domain of abstract interpretation~\cite{cousot77}, retains
them: multiplication of signs is exact, and the only sum whose sign is not determined is $+$ plus $-$, which
is $?$. Taking signs as the scalars gives the \emph{qualitative derivatives} of qualitative
reasoning~\cite{dekleer84,kuipers86}, an entry reading ``increasing this input increases (or decreases, or
cannot affect) that output'', with $?$ appearing exactly at the potential cancellations: in the refund
example, the entry for the price becomes $?$. Collapsing $+$, $-$ and $?$ to $1$ recovers the Boolean
analysis.

\subsubsection{Linearised provenance}
\label{sec:approx-as-tangents:linearised-provenance}

Database provenance has its own semiring tradition: in the framework of
\citet{green07}, query results carry annotations from a commutative semiring, with the \emph{free}
commutative semiring $\mathbbm{N}[X]$ of \emph{provenance polynomials} playing a universal role: its
variables name the input tuples, multiplication records their joint use in a derivation, and addition
collects the alternative derivations. By freeness, assigning a value to each variable extends uniquely to a
semiring homomorphism out of $\mathbbm{N}[X]$, so provenance computed once as polynomials determines, by
evaluation, the annotation in every other semiring.

Provenance semirings can serve as the scalars of the present framework, but with a limitation: a
derivative is linear, so the framework captures \emph{linearised} provenance. Taking $S = \mathbbm{N}[X]$, a
Jacobian entry is a polynomial recording symbolically how an input position contributes to an output position.
Green's semantics assigns an output a polynomial $p$ whose degree records joint use; seeding each input
position with a unique variable, our forward analysis instead returns a polynomial linear in the position
variables, with the partial derivative $\partial p / \partial x_i$, evaluated at the run's input values, as
the coefficient of $x_i$. Joint use survives only in the coefficients: an input
used twice contributes $2x$ where the annotation would record $x^2$.

Translation along a homomorphism applies to the present framework as well: a semiring homomorphism acts
entrywise on Jacobians, and the action is functorial, since matrix composition is built from addition and
multiplication. The translation of interest for dependency tracking is ``zero-testing'' into the Booleans, and
whether that is a homomorphism depends on the source semiring. The counting semiring $\mathbbm{N}$, itself a
provenance semiring recording how many derivations use each input, is \emph{positive}: a sum is zero only when
both summands are, and a product only when a factor is; equivalently, zero-testing is a homomorphism, so
counting translates exactly to dependency tracking. The rationals are not positive, since a number and its
negation sum to zero, and zero-testing $\mathbbm{Q} \to \Two$ preserves multiplication exactly but addition
only laxly. It still acts entrywise on matrices, preserving identities and transposition, but composition only
up to over-approximation: composing Boolean Jacobians can record dependencies absent from the Boolean image of
their composite. The refund in \exref{introduction-example} is a case in point: the price genuinely does not
matter to the total, yet the composed dependencies say it might.

\section{Models of Galois Slicing for a Total Language} % \GPS doesn't use title caps
\label{sec:models-of-total-gps}

\Replaced{\secref{dependency-as-differentiation} presented the dependency maps of \GPS as conjugate pairs of
join-preserving maps, represented by Boolean matrices. The category $\Mat(\Two)$ is not itself a suitable
setting for interpreting a programming language: it has biproducts, so products and coproducts coincide,
presenting difficulties for interpreting a language with coproducts or exponentials. Moreover, just as a Jacobian describes a smooth
function only at a point, a Boolean matrix captures the dependency structure of a program only at a single
input value; a model must package the underlying function on values together with a matrix at each input.}{The previous section concluded that L-posets and stable functions give
a model of \GPS analogous to manifolds and smooth functions. However,
we noted a conceptual shortcoming of this model, for the purposes of
modelling total computations, that proper values and their
approximations live in the same category.} In this section, we propose
a model for total \GPS based on the \emph{Category of Families}
construction. This construction, and the more general Grothendieck
construction, has been previously used by V{\'a}k{\'a}r and
collaborators \cite{vakar22} to model automatic differentiation for
higher-order programs on the reals. We reuse some of their results,
and discuss the commonalities as we go.

\subsection{The Category of Families Construction}
\label{sec:models-of-total-gps:fam}

\Replaced{Our model pairs a set $X$ of values with, for each $x \in X$, a bounded lattice $\partial X(x)$ of
approximations of $x$, generalising the lattices of positions on which a Boolean Jacobian acts in
\secref{dependency-as-differentiation}.}{L-posets are partially ordered sets where every principal downset
$\downset{x}$ is a bounded lattice of approximations/tangents. As we
explained in \secref{stab}, the shortcoming of this setup is that
proper values and their approximations live in the same set. We fix
this by changing our model to one where we have sets $X$ of values,
and for each $x \in X$, a bounded lattice $\partial X(x)$ of
approximations of $x$.} This construction is an instance of the general
\emph{Category of Families} construction:

\begin{definition}
  Let $\cat{C}$ be a category. The \emph{Category of Families} over
  $\cat{C}$, $\Fam(\cat{C})$, has as objects pairs $(X, \partial X)$,
  where $X$ is a set and $\partial X : X \to \cat{C}$ is an
  $X$-indexed family of objects in $\cat{C}$. A morphism
  $f : (X, \partial X) \to (Y, \partial Y)$ consists of a pair of a
  function $f : X \to Y$ and a family of morphisms of $\cat{C}$,
  $\partial f : \Pi_{x : X}.\,\cat{C}(\partial X(x), \partial Y(f\,x))$.
\end{definition}

The reason for choosing the $\Fam$ construction is that composition in
this category is an abstract version of the chain rule that we have
seen in \Replaced{\remref{chain-rule} and \remref{chain-rule-bool}}{\remref{chain-rule},
\remref{chain-rule-cm}, and \remref{chain-rule-stable}}. Composition $f \circ g $ of morphisms
$f : (Y, \partial Y) \to (Z, \partial Z)$ and
$g : (X, \partial X) \to (Y, \partial Y)$ in this category is given by
normal function composition on the set components, and
$\partial (f \circ g)(x) = \partial f(f, x) \circ \partial g(x)$,
where the latter composition is in $\cat{C}$.

The fact that morphisms in $\Fam(C)$ compose according to a chain rule
means that the categories we considered in \secref{approx-as-tangents}
embed into $\Fam(C)$ for appropriate $C$. If we let $\FinVect$ be the
category of finite dimensional real vector spaces and linear maps,
then:

\begin{proposition}
  \label{prop:embed-manifolds}
  There is a faithful functor $\Man \to \Fam(\FinVect)$ that sends a
  manifold $M$ to $(M, \lambda x. T_x(M))$, and each smooth function
  $f$ to $(f, f_*)$, the pair of $f$ and its forward derivative.
\end{proposition}

A similar result is given by \citet{cruttwell2022}, where Euclidean
spaces $\RR^n$ and smooth functions are embedded into a category of
lenses (the ``simply typed'' version of the $\Fam$ construction). As
in \citet{vakar22}, the idea is to formally separate functions on
points and their forward/reverse tangent maps for the purposes of
implementation of automatic differentiation. In the case of smooth
maps, this process throws away information on higher derivatives by
turning smooth maps into pairs of plain functions and linear
functions. We conjecture at the end of this section that the analogous
construction in our partially ordered setting does not.

\Replaced{For dependency analysis, we pick categories of lattices and join-preserving maps:}{For the
categories $\CM$ and $\Lposet$, we pick the appropriate categories of partial orders and monotone maps:}

\begin{enumerate}[leftmargin=\enummargin]
\item \Replaced{The category $\LatConj$ has bounded lattices as objects and conjugate pairs
  (\defref{conjugate-pair}) as morphisms, with the forward map going in the direction of the morphism. The
  category $\Fam(\LatConj)$ is our preferred model for total functions with dependency analysis. We explore
  some specific examples in this category in \secref{semantic-gps}.}{The category $\LatGal$ has bounded
  lattices as objects and Galois connections as morphisms, with the right adjoint going in the ``forward''
  direction. The category $\Fam(\LatGal)$ is our preferred model for total functions with Galois slicing. We
  explore some specific examples in this category in \secref{semantic-gps}.}
\item \Replaced{The category $\JoinSLat$ has join-semilattices with bottom as objects and monotone finite
  join preserving functions as morphisms. The categories $\Fam(\JoinSLat)$ and $\Fam(\JoinSLat^\op)$ provide
  models of total functions with only ``forward derivatives'' or only ``backwards derivatives'',
  respectively.}{The category $\MeetSLat$ has meet-semilattices with top as objects and monotone finite meet
  preserving functions as morphisms. The category $\Fam(\MeetSLat)$ provides a model of total functions with
  ``forward derivatives'' only. The category $\JoinSLat$ has join-semilattices with bottom as objects and
  monotone finite join preserving functions as morphisms. The category $\Fam(\JoinSLat^\op)$ provides a model
  of total functions with ``backwards derivatives'' only.}
\end{enumerate}

\Deleted{We can get an analogous result to \propref{embed-manifolds} for L-posets and stable maps: there is
a faithful functor $\Lposet \to \Fam(\LatGal)$ that maps an L-poset $X$ to $(X, \lambda x. \downset{x})$ and
stable functions $f$ to $(f, \lambda x. (f_x, f^*_x))$, and likewise a faithful functor
$\CM \to \Fam(\MeetSLat)$. Despite its similarity, this result has a lesser status than
\propref{embed-manifolds} because it is not clear that the category $\Lposet$ (or $\CM$) is a canonical
definition of approximable sets and functions with approximation derivatives, as we discussed at the end of
\secref{stab}.}
Our working hypothesis is that $\Fam(\Replaced{\LatConj}{\LatGal})$, where values and their
approximations are separated by construction, is a natural model of
semantic \GPS in a total setting, though we note some shortcomings in
\remref{further-structure}. We now investigate some categorical
properties of this category, with a view to modelling a higher-order
total programming language in \secref{language}.

\subsection{Categorical Properties of $\Fam(\cat{C})$}

\subsubsection{Coproducts and Products}
\label{sec:models-of-total-gps:coproducts-and-products}

The categories $\Fam(\cat{C})$ are the free coproduct completions of
categories $\cat{C}$, so they have all coproducts:

\begin{proposition}
  For any $\cat{C}$, $\Fam(\cat{C})$ has all coproducts, which can be
  given on objects by:
  \begin{displaymath}
    \coprod_i (X_i, \partial X_i) = (\coprod_i X_i, \lambda (i, x_i).\, \partial X_i(x))
  \end{displaymath}
  Coproducts in $\Fam(\cat{C})$ are \emph{extensive}
  \cite[Proposition 2.4]{carboni-lack-walters93}.
\end{proposition}

For $\Fam(\cat{C})$ to have finite products, we need $\cat{C}$ to have
finite products:

\begin{proposition}
  If $\cat{C}$ has finite products, then so does $\Fam(\cat{C})$. On
  objects, binary products can be defined by:
  \begin{displaymath}
    (X, \partial X) \times (Y, \partial Y) = (X \times Y, \lambda (x, y). \partial X(x) \times \partial Y(y))
  \end{displaymath}
  Since $\Fam(\cat{C})$ is extensive, products and coproducts
  distribute.
\end{proposition}

Using the infinitary coproducts and finite products, we can construct
a wide range of other useful semantic models of datatypes in
$\Fam(\cat{C})$. For example, lists can be constructed as a coproduct
\begin{equation}
  \label{eqn:lists}
  \List(X) = \coprod_{n \in \mathbb{N}} X^n
\end{equation}
where $X^0 = 1$ (the terminal object) and $X^{n+1} = X \times X^n$.

\Replaced{Our category of interest, $\Fam(\LatConj)$ has all coproducts, from the $\Fam$ construction
itself, and finite products, because $\LatConj$ has products. Similarly for $\Fam(\JoinSLat)$. As we shall
see below, the products in $\LatConj$ (and
$\JoinSLat$) are also coproducts, which is essential to obtaining Cartesian closure.}{Our category of
interest, $\Fam(\LatGal)$ has coproducts and finite products, because $\LatGal$ has products. Similarly for
$\Fam(\MeetSLat)$. As we shall see below, the products in $\LatGal$ (and $\MeetSLat$ and $\JoinSLat)$ are
also coproducts, which is essential to obtaining Cartesian closure.}

\subsubsection{Cartesian Closure}

For Cartesian closure of the categories $\Fam(C)$ that we are
interested in, we rely on the following theorem of \citet{nunes2023},
specialised from their setting with the general Grothendieck
construction to $\Fam(\cat{C})$. This relies on the definition of
\emph{biproducts}, which we discuss below in \secref{biproducts}.

\begin{theorem}[\cite{nunes2023}]
  \label{thm:fam-closed}
  \AGDA.  If $\cat{C}$ has biproducts (\defref{biproducts}) and all
  products, then $\Fam(\cat{C})$ is Cartesian closed\footnote{More
    precisely, if $\cat{C}$ has coproducts then we have a monoidal
    product on $\Fam(\cat{C})$ which is closed by this
    construction. When these coproducts are in fact biproducts, we get
    Cartesian closure.}. On objects, the internal hom can be given by:
  \begin{displaymath}
    (X, \partial X) \to (Y, \partial Y) = (\Pi_{x : X}. \Sigma_{y : Y}. \cat{C}(\partial X(x), \partial Y(y)), \lambda f. \Pi_{x : X}. \partial Y(\pi_1(f\, x)))
  \end{displaymath}
\end{theorem}

The $\Set$-component of $(X, \partial X) \to (Y, \partial Y)$ consists
of exactly the morphisms of $\Fam(\cat{C})$, rephrased into a single
object. When $\cat{C} = \FinVect$, these are functions with an
associated linear map at every point, and when
\Replaced{$\cat{C} = \LatConj$, these are functions with an associated conjugate pair at every
point}{$\cat{C} = \LatGal$, these are functions with an associated Galois connection at every point}. A tangent to a function is then defined to be a mapping from
points in the domain to tangents in the codomain along the function.

The category \Replaced{$\JoinSLat$}{$\MeetSLat$} satisfies the hypotheses of
\thmref{fam-closed}, so:
\begin{corollary}
  $\Fam(\Replaced{\JoinSLat}{\MeetSLat})$ is Cartesian closed and has all coproducts.
\end{corollary}
Unfortunately, neither \Replaced{$\LatConj$}{$\LatGal$} nor $\FinVect$ satisfy the
hypotheses of this theorem, because neither of them have infinite
products. We will consider ways to rectify this below in
\secref{fixing-completeness}.

\begin{remark}
  \label{rem:hermida-exponentials}
  There is another construction of internal homs on $\Fam(\cat{C})$
  arising from the use of fibrations for categorical logical
  relations, due to \citet[Corollary 4.12]{hermida99}. If we assume
  that $\cat{C}$ is itself Cartesian closed and has all products, then
  we could construct an internal hom as:
  \begin{displaymath}
    (X, \partial X) \to (Y, \partial Y) = (X \to Y, \lambda f. \Pi_{x : X}.\,\partial X(x) \to \partial Y(f\,x))
  \end{displaymath}
  However, for the purposes of modelling differentiable programs, this
  is fatally flawed in that neither \Replaced{$\LatConj$}{$\LatGal$} nor $\FinVect$ are
  Cartesian closed, and there is no way of making them so without
  losing the property of being able to conjunct or add tangents, as we
  shall see below. We will implicitly use Hermida's construction in
  our definability proof in \secref{definability}, where we use a
  logical relations argument to show that every morphism definable in
  the higher order language is also first-order definable.
\end{remark}

\subsubsection{CMon-Categories and Biproducts}
\label{sec:biproducts}

Loosely stated, biproducts are objects that are both products and
coproducts. The concept can be defined in any category, as shown by
\citet{karvonen20}, but for our purposes it will be more convenient to
use the shorter definition in categories enriched in commutative
monoids:
\begin{definition}
  A category $\cat{C}$ is enriched in $\CMon$, the category of
  commutative monoids, if every homset $\cat{C}(X,Y)$ is a commutative
  monoid with $(+,0)$ and composition is bilinear:
  \begin{displaymath}
    f \comp \zero = 0 = \zero \comp f
  \end{displaymath}
  \begin{displaymath}
    (f + g) \comp h = (f \comp h) + (g \comp h) \qquad
    h \comp (f + g) = (h \comp f) + (h \comp g)
  \end{displaymath}
\end{definition}
In any $\CMon$-category we can define what it means to be the
biproduct of two objects:
\begin{definition}
  \label{def:biproducts}
  In a $\CMon$-category a biproduct is an object $X \biprod Y$
  together with morphisms

  \begin{center}
    \begin{tikzcd}
      X \arrow[r, "\biinj_X", shift left] &
      X \biprod Y \arrow[l, "\biproj_X", shift left] \arrow[r, "\biproj_Y"', shift right] &
      Y \arrow[l, "\biinj_Y"', shift right]
    \end{tikzcd}
  \end{center}

  \vspace{-1mm}
  \noindent satisfying

  \vspace{-5mm}
  \begin{minipage}[t]{0.45\textwidth}
    \begin{center}
      \begin{salign*}
        \biproj_X \comp \biinj_X &= \id_X \\
        \biproj_Y \comp \biinj_X &= \zero_{X,Y}
      \end{salign*}
    \end{center}
  \end{minipage}%
  \begin{minipage}[t]{0.45\textwidth}
    \begin{center}
      \begin{salign*}
        \biproj_Y \comp \biinj_Y &= \id_Y \\
        \biproj_X \comp \biinj_Y &= \zero_{Y,X}
      \end{salign*}
    \end{center}
  \end{minipage}

  \begin{salign*}
    (\biinj_X \comp \biproj_X) + (\biinj_Y \comp \biproj_Y) &= \id_{X \biprod Y}
  \end{salign*}
  A zero object is an object that is both initial and terminal.
\end{definition}
As the name suggests, biproducts in a category are both products and
coproducts:
\begin{proposition}
  \item
  \begin{enumerate}[leftmargin=\enummargin]
  \item A $\CMon$-category that has biproducts $X \oplus Y$ for all
    $X$ and $Y$ also has products and coproducts with
    $X \times Y = X + Y = X \oplus Y$.
  \item A $\CMon$-category with (co)products also has biproducts, and
    any initial or terminal object is a zero object.
  \end{enumerate}
\end{proposition}

\begin{example}
  The following are $\CMon$-enriched and have finite products, and
  hence biproducts:
  \begin{enumerate}[leftmargin=\enummargin]
  \item In $\FinVect$, morphisms are linear maps and so can be added
    and have a zero map. Finite products are given by Cartesian
    products of the underlying sets, with the vector operations
    defined pointwise.
  \item \Replaced{In $\LatConj$, both components of a conjugate pair are summed using joins, and the zero
    maps are given by the constantly $\bot$ functions.}{In $\LatGal$, right
    adjoints are summed using meets and left adjoints are summed using joins. The zero maps are given by the
    constantly $\top$ and constantly $\bot$ functions respectively.} Products are given by the Cartesian
    product of the underlying set and the one-element lattice for the terminal/initial/zero object.
  \item \Replaced{$\JoinSLat$ is also $\CMon$-enriched and has finite products similar to
    $\LatConj$.}{$\MeetSLat$ and $\JoinSLat$ are both $\CMon$-enriched and have finite products similar to
    $\LatGal$.}
  \end{enumerate}
\end{example}

\begin{remark}
  Categories with zero objects cannot be Cartesian closed without
  being trivial in the sense of having exactly one morphism between
  every pair of objects because
  $\cat{C}(X, Y) \cong \cat{C}(1 \times X,Y) \cong \cat{C}(0 \times
  X,Y) \cong \cat{C}(0,X \to Y) \cong 1$. Consequently, we cannot
  apply the alternative construction of exponentials described in
  \remref{hermida-exponentials}.
\end{remark}

\subsubsection{Discrete Completeness}
\label{sec:fixing-completeness}

The second hypothesis of \thmref{fam-closed} is that the category
$\cat{C}$ has all (i.e., infinite) products. This is required to
gather together tangents for all of the points in the domain of the
function. Unfortunately, neither $\FinVect$ nor \Replaced{$\LatConj$}{$\LatGal$} is complete
in this sense.

In the case of $\FinVect$, the solution is to expand to the category
of all vector spaces $\Vect$, where infinite direct products
exist. Note that these infinite products are not biproducts because
the vector space operations themselves are finitary. This is the
solution that \citet{vakar22} use for the semantics of forward
($\Fam(\Vect)$) and reverse ($\Fam(\Vect^\op)$) automatic
differentiation for higher order programs. Since the forward and
reverse derivatives of a smooth map are intrinsically defined,
\citet{vakar22}'s correctness theorem shows that, for programs with
first-order type, the interpretation in $\Fam(\Vect)$ correctly yields
the forward derivative of the defined function on the reals (and
reverse derivative for $\Fam(\Vect^\op)$).

For \Replaced{$\LatConj$}{$\LatGal$} we could expand to the category of complete lattices and
\Replaced{conjugate pairs}{Galois connections} between them. From a classical mathematical point
of view, this would give a model of \GPS that would be suitable for
reasoning about programs' behaviour and their forward and backward
approximations. However, in terms of building an executable model
inside the Agda proof assistant, and with an eye toward implementation
strategies, we seek a finitary solution. (Note that the solution of
moving to complete lattices is very different to moving to arbitrary
dimension vector spaces: in the former we have infinitary operations,
while the latter still has only finitary operations.)

We will avoid the need for infinitary operations by separating the forward and backward parts of
the \Replaced{conjugate pairs}{Galois connections} to act independently by moving to the product category
\Replaced{$\JoinSLat \times \JoinSLat^\op$}{$\MeetSLat \times \JoinSLat^\op$}. Objects in this category
consist of \emph{separate} \Replaced{join-semilattices and potentially unrelated forward and backward
join-preserving maps}{meet- and join-semilattices and potentially unrelated forward meet-preserving and
backward join-preserving maps}. We first
check that this category satisfies the hypotheses of
\thmref{fam-closed}:

\begin{proposition}
  \Replaced{$\JoinSLat \times \JoinSLat^\op$}{$\MeetSLat \times \JoinSLat^\op$} has biproducts and all products.
\end{proposition}

\begin{proof}
  \Replaced{$\JoinSLat$ is $\CMon$-enriched and has}{$\MeetSLat$ and $\JoinSLat$ are both $\CMon$-enriched
  and have} finite products, as noted above. The opposite of a category with biproducts also has biproducts
  (by swapping the injections $i$ and projections $p$), and products of categories with biproducts also have
  biproducts pointwise. Hence \Replaced{$\JoinSLat \times \JoinSLat^\op$}{$\MeetSLat \times \JoinSLat^\op$}
  has biproducts.

  \Replaced{$\JoinSLat$}{$\MeetSLat$} has all products, indeed all limits, because it is the category of
  algebras for a Lawvere theory. \Replaced{It also has}{Similarly, $\JoinSLat$ has} all coproducts, indeed
  all colimits, for the same reason. Note that these are very different constructions: elements of a product
  of \Replaced{join-semilattices}{meet-semilattices} consist of (possibly infinite) tuples of elements,
  while elements of a coproduct of join-semilattices consist of \emph{finite} formal joins of elements
  quotiented by the join-semilattice equations. Since $\JoinSLat$ has all coproducts, $\JoinSLat^\op$ has
  all products, and so \Replaced{$\JoinSLat \times \JoinSLat^\op$}{$\MeetSLat \times \JoinSLat^\op$} has all
  products, as required.
\end{proof}

\begin{corollary}
  \label{cor:mslat-jslat-bcc}
  $\Fam(\Replaced{\JoinSLat}{\MeetSLat} \times \JoinSLat^\op)$ is Cartesian closed and has
  all coproducts.
\end{corollary}

This corollary means that, assuming a sensible intepretation of primitive types and operations, we
can use $\Fam(\Replaced{\JoinSLat}{\MeetSLat} \times \JoinSLat^\op)$ to
interpret the higher-order language we describe in the next section. We still regard the category
$\Fam(\Replaced{\LatConj}{\LatGal})$ as the reference model of approximable sets with
forward and backward approximation maps; the category
$\Fam(\Replaced{\JoinSLat}{\MeetSLat} \times \JoinSLat^\op)$ is a technical
device to carry out the interpretation of higher-order programs. To get interpretations of first-order types
and primitive operations, we can embed $\Fam(\Replaced{\LatConj}{\LatGal})$ into
$\Fam(\Replaced{\JoinSLat}{\MeetSLat} \times \JoinSLat^\op)$:

\begin{proposition}
  \label{prop:ho-embedding}
  The functor
  $\HoEmbed : \Fam(\Replaced{\LatConj}{\LatGal}) \to \Fam(\Replaced{\JoinSLat}{\MeetSLat} \times
  \JoinSLat^\op)$ is defined
  on objects as
  $\HoEmbed(X, \partial X) = (X, \lambda x. (\partial X(x), \partial
  X(x)))$. This functor is faithful and preserves coproducts and
  finite products.
\end{proposition}

With this embedding functor, we will see in \lemref{first-order-agreement-types} that the
interpretation of first-order types will be the same up to isomorphism in
$\Fam(\Replaced{\LatConj}{\LatGal})$ and
$\Fam(\Replaced{\JoinSLat}{\MeetSLat} \times \JoinSLat^\op)$, as long as we
interpret the base types as objects in $\Fam(\Replaced{\LatConj}{\LatGal})$. At higher-order,
however, the \Replaced{forward and backward}{meet-semilattice and join-semilattice} sides of the
interpretation will diverge, and it is no longer clear that the interpretation of programs using
higher-order functions internally will result in \Replaced{conjugate pairs}{Galois connections}. In
\secref{definability} we will see that every program with first-order type (even if it uses higher-order
functions internally) does in fact have an interpretation definable in
$\Fam(\Replaced{\LatConj}{\LatGal})$.

% \bob{In the conclusion, we have to admit to the possibility that
%   $\Fam(CLatGal)$ might be sensible. We could have also used
%   PShCMon(LatGal) and there is PShCMon(GraphLang) for
%   normalisation. We would have different problems with proving
%   correctness however.}

\subsection{Semantic Galois Slicing in $\Fam(\Replaced{\LatConj}{\LatGal})$}
\label{sec:semantic-gps}

Our thesis is that $\Fam(\Replaced{\LatConj}{\LatGal})$ is a suitable setting for
interpreting first-order programs for \GPS. The above discussion has
been somewhat abstract, so we now consider some examples in the
category $\Fam(\Replaced{\LatConj}{\LatGal})$ and how they relate to \GPS.

Spelt out in full, $\Fam(\Replaced{\LatConj}{\LatGal})$ has as objects $(X, \partial X)$,
all pairs of a set $X$ and and for every $x \in X$, a bounded lattice
$\partial X(x)$. Morphisms $(X, \partial X) \to (Y, \partial Y)$, are
triples $(f, \partial f_f, \partial f_r)$ of functions $f : X \to Y$
and families of \Replaced{join-preserving}{monotone} maps
$\partial f_f : \Pi_{x : X}.\partial X(x) \multimap \partial Y(f\,x)$
(``forward derivative'') and
$\partial f_r : \Pi_{x : X}. \partial Y(f\,x) \multimap \partial X(x)$
(``reverse derivative''), \Replaced{such that for all $x$, $\partial f_f(x)$ and $\partial f_r(x)$ form a
conjugate pair}{such that for all $x$, $\partial f_r(x) \dashv \partial f_f(x)$}.

\subsubsection{Unapproximated Functions}
\label{sec:unapproximated-functions}

The category \Replaced{$\LatConj$}{$\LatGal$} has a terminal (also zero) object $\mathbb{1}$, so there is
a functor $\Disc : \Set \to \Fam(\Replaced{\LatConj}{\LatGal})$ that maps a set $X$ to
$(X, \lambda x. \mathbb{1})$ and functions $f$ to morphisms
$(f, \lambda\_.\,\mathrm{id}_{\mathbb{1}})$. This functor preserves
products and coproducts. Therefore, we can take any sets and functions
of interest for modelling primitive types and operations of a
programming language and embed, albeit without any interesting
approximation information.

\subsubsection{Lifting Monad}
\label{sec:models-of-total-gps:lifting}

\Added{\todo{Revisit this. There doesn't seem to be a lifting monad in  $\LatConj$; ``add new $\top$'' has a unit but isn't
functorial, and ``add a new $\bot$'' is a comonad, not a monad (it has no unit). So we can't express these slicing behaviours in the
language, using lifting; they would have to be given the desired short-circuiting dependency behaviour directly.}}

The operation of adding a new bottom element to a bounded lattice
forms part of a monad $L$ on $\LatGal$. This monad extends to a
(strong) monad $\Lift$ on $\Fam(\LatGal)$ with
$\Lift(X, \partial X) = (X, L \circ \partial X)$. The monad $\Lift$
does not affect the points of the original object, but adds a new
minimum approximation.

\newcommand{\Bool}{\mathrm{Bool}}

Let $\Bool = \Disc(\{\mathsf{tt},\mathsf{ff}\})$ be the
(unapproximated) embedding of the Booleans and
$\mathrm{or} : \Bool \times \Bool \to \Bool$ be the (unapproximated)
Boolean OR function. Using a Moggi-style let notation
\cite{notions-of-computation} for morphisms constructed using the
Monad structure of $\Lift$, we can reproduce the functions
$\mathrm{strictOr}$ and $\mathrm{shortCircuitOr}$ functions from
\exref{strict-short-circuit} (we also assume an if-then-else operation
on Booleans, definable from the fact that $\Fam(\LatGal)$ has
coproducts and $\Disc$ preserves them). Both of these expressions
define morphisms $\Lift(\Bool) \times \Lift(\Bool) \to \Lift(\Bool)$
in $\Fam(\LatGal)$:
\begin{displaymath}
  \begin{array}{lcl}
    \mathrm{strictOr}(x,y)&=&\mathrm{let}\,b_1 \Leftarrow x\,\mathrm{in}\,\mathrm{let}\,b_2 \Leftarrow y\,\mathrm{in}\,\eta(\mathrm{or}(b_1,b_2)) \\
    \mathrm{shortCircuitOr}(x,y)&=&\mathrm{let}\,b_1 \Leftarrow x\,\mathrm{in}\,\mathrm{if}\,b_1\,\mathrm{then}\,\eta(\mathsf{tt})\,\mathrm{else}\,y
  \end{array}
\end{displaymath}
Examining the morphisms so defined in $\Fam(\LatGal)$, we can see
that, in the $\Set$ component, they are both exactly the normal
Boolean-or operation. However, they have different approximation
behaviour, reflecting the different ways that they examine their
inputs. Let us write $\top,\bot$ for the elements of the approximation
lattice at each point of $\Lift(\Bool)$, then applying the reverse
derivative at $(\mathsf{tt},\mathsf{tt})$ to the tangent $\top$
reveals which of the inputs contributed to the output for each
function:
\begin{displaymath}
  \begin{array}{lcl}
    (\partial \mathrm{strictOr})_r(\mathsf{tt},\mathsf{tt})(\top) &=& (\top, \top)  \\
    (\partial \mathrm{shortCircuitOr})_r(\mathsf{tt},\mathsf{tt})(\top) &=& (\top, \bot)
  \end{array}
\end{displaymath}
In comparison to the categories $\CM$ and $\Lposet$ from
\secref{approx-as-tangents}, we have retained the usage information in
the forward and reverse tangents, but we also accurately model
totality of the functions. That is, the constantly $\bot$ function is
also present in both $\CM$ and $\Lposet$, but is not expressible in
$\Fam(\LatGal)$.

An analogue of the $\mathrm{parallelOr}$ function from
\exref{parallel-or} is not definable in $\Fam(\LatGal)$. We would have
to have
$(\partial \mathrm{parallelOr})_f(\mathsf{tt},\mathsf{tt})(\top,\bot)
= (\partial \mathrm{parallelOr})_f(\mathsf{tt},\mathsf{tt})(\bot,\top)
= \top$ to reflect the desired property that either of the inputs
being $\mathsf{tt}$ is enough to determine the output. We also must
have
$(\partial \mathrm{parallelOr})_f(\mathsf{tt},\mathsf{tt})(\bot,\bot)
= \bot$, to reflect the fact that we will get no information in the
output if we required that neither of the inputs is examined. However,
this means that
$(\partial \mathrm{parallelOr})_f(\mathsf{tt},\mathsf{tt})$ will not
preserve meets because $(\top,\bot) \sqcap (\bot,\top) = (\bot,\bot)$
but $\top \neq \bot$.

An analogue of the $\mathrm{gustave}$ function from
\exref{parallel-or} is definable in $\Fam(\LatGal)$, but not using the
lifting monad structure as we could for $\mathrm{strictOr}$ and
$\mathrm{shortCircuitOr}$.

\begin{remark}
  \label{rem:further-structure}
  These examples highlight a potential criticism of $\Fam(\LatGal)$ as
  a category for modelling \GPS. For $\mathrm{shortCircuitOr}$, we had
  $(\partial \mathrm{shortCircuitOr})_r(\mathsf{tt},\mathsf{tt})(\top)
  = (\top, \bot)$, indicating that the second argument was not needed
  for computing the output. However, there is no way, in
  $\Fam(\LatGal)$, of turning this into a rigorous statement that the
  $\Set$-component of this morphism does not actually depend on its
  second argument. We conjecture that this can be rectified by
  requiring some kind of additional structure on each object
  $(X, \partial X)$ consisting of a map
  $\Pi_{x : X}.\partial X(x) \to \mathcal{P}(X)$, where $\mathcal{P}$
  is the powerset, which identifies for each $x$ which elements are
  indistinguishable from $x$ at this level of approximation. One would
  also presumably have to require additional conditions for this to
  respect the lattice structure and be preserved by
  morphisms\footnote{This additional structure is reminiscent of the
    additional structure on \emph{directed containers} defined by
    \citet{ahman-chapman-uustalu2012}. They require a map
    $\Pi_{x : X}.\partial X(x) \to X$ picking out a specific $X$
    ``jumped to'' by some change $\delta x : \partial X(x)$ at $x$. In
    our proposed setup, we follow the approximation theme of \GPS by
    having a \emph{set} of things that could be used to replace the
    original $x$. We observe that objects of $\Fam(\LatGal)$ arising
    from $\Lposet$ are ``directed'' in the
    \citet{ahman-chapman-uustalu2012} sense because the map can pick
    out the element of the original poset that was approximating
    $x$.}.
\end{remark}

The $\Lift$ monad provides a controllable way of adding
presence/absence approximation points to composite data, and its monad
structure makes explicit in the program structure exactly how such
approximations are propagated through computations. The fact that
there are different choices of this kind of approximation tracking
provides freedom to the language implementor to decide what
information is worth tracking. The Galois slicing implementations
discussed in \citet{perera12a} and \citet{ricciotti17}, for example,
bake-in an approximation point at every composite type constructor. We
will see in \secref{cbn-translation} that this choice can be
systematised in our setting by considering a monadic CBN translation
to uniformly add $\Lift$ approximation points to composite data types.

\subsubsection{An Approximation Object and the Tagging Monad}
\label{sec:tagging-monad}

\Replaced{We take as given an object $\mathbb{A} = (1, \lambda\_.\,\{\top,\bot\})$ for tagging first-order
data with presence and absence information.}{The $\Lift$ monad provides a way of tagging first-order data with
presence and absence information. The object $\Lift(1)$, the lifting
of the terminal object in $\Fam(\LatGal)$, yields an object that
consists purely of presence/absence information:
$\mathbb{A} = (1, \lambda\_.\,\{\top,\bot\})$.} This object is the
carrier of a commutative monoid in $\Fam(\Replaced{\LatConj}{\LatGal})$, where the forward
maps take the \Replaced{join (the output is present if either input is)}{meet (both the inputs are required for the output to be
present)} and the backwards maps duplicate.% \bob{Is it also a comonoid?}

Since $\mathbb{A}$ is a monoid, we can define the writer monad
$T(X) = \mathbb{A} \times X$ in $\Fam(\Replaced{\LatConj}{\LatGal})$ which ``tags'' $X$
values with approximation information.
\Added{\todo{Revisit the comparison with the lifting monad $L$ below (see \secref{models-of-total-gps:lifting}).}}
This is similar to the $\Lift$
monad in that it adds approximation information to an object. On
discrete objects it agrees with the lifting:
$T(\Disc(X)) \cong L(\Disc(X))$. However, on composite data, the two
monads give different approximation lattices. Let $A$ and $B$ be sets.
Then $T(T(\Disc(A)) \times T(\Disc(B)))$ has approximation lattices
at $(a,b)$ that are always isomorphic to $\{\top,\bot\}^3$. The
corresponding $\Lift(\Lift(\Disc(A)) \times \Lift(\Disc(B)))$ object's
approximation lattice at $(a,b)$ is always isomorphic to
$(\{\top,\bot\}^2)_\bot$.
In terms of usage tracking, the object using the $\Lift$ monad is more
appealing. The approximation lattice resulting from the use of the $T$
monad contains apparently nonsensical elements corresponding to
``using'' one or other components of the product \emph{without using
  the product itself}. (This arises from the fact that we can project
the $X$ out of $\mathbb{A} \times X$ without touching the
$\mathbb{A}$.)


This example shows that we have to be careful about how we choose the
interpretation of approximable sets in $\Fam(\Replaced{\LatConj}{\LatGal})$, and again
highlights the point we made in \remref{further-structure} that
perhaps $\Fam(\Replaced{\LatConj}{\LatGal})$ does not have quite enough structure to
determine ``sensible'' approximation information. On the other hand,
the use of the $T$ monad does have the advantage that the
approximation lattices built from discrete sets, products, and
coproducts, are always Boolean lattices, meaning that we can take
complements of usage information. The ability to take complements of
approximations has been used by \citet{perera22} to compute
\emph{related} outputs, as we discuss in
\secref{related-work:galois-slicing}.

\subsubsection{Approximating Numbers by Intervals}
\label{sec:interval-approx}

So far, the approximation lattices we have looked at in
$\Fam(\Replaced{\LatConj}{\LatGal})$ have only consisted of those constructed from finite
products and lifting, and only track binary usage/non-usage
information. \exref{intervals-and-maxima-elements} shows how we can go
beyond this to get more ``quantitative'' approximation
information. Let the object of reals with interval approximations in
$\Fam(\Replaced{\LatConj}{\LatGal})$ be
$\mathbb{R}_{\mathit{intv}} = (\mathbb{R}, \lambda x.\,\{[l,u] \mid l
\leq x \leq u\} \cup \{\bot\})$ where the lattices of intervals are
reverse ordered by inclusion with $\bot$ at the bottom. Then,
following the examples in \exref{intervals-and-maxima-elements} and
\exref{stable-functions}, we can define addition, negation, and
scaling by $r \geq 0$:
\cbstart
\begin{displaymath}
  \begin{array}{lcl}
    \mathrm{add}&=&(\begin{array}[t]{@{}l}
      \lambda (x_1,x_2).\,x_1+x_2, \\
      \lambda (x_1,x_2)\,([l_1,u_1],[l_2,u_2]).\,[(l_1+x_2) \sqcup (l_2+x_1),(u_1+x_2) \sqcap (u_2+x_1)] \\
      \lambda (x_1,x_2)\,[l,u].\,([l-x_2,u-x_2],[l-x_1,u-x_1]))
    \end{array}\\
    \mathrm{neg}&=&(\lambda x.\,{-}x, \lambda x\,[l,u].\,[-u,-l], \lambda x\,[l,u].\,[-u,-l]) \\
    \mathrm{scale}(r)&=&(\lambda x.\,rx, \lambda x\,[l,u].\,[rl,ru], \lambda x\,[l,u].\,\mathrm{if}\,r=0\,\mathrm{then}\,\bot\,\mathrm{else}\,[\frac{l}{r},\frac{u}{r}])
  \end{array}
\end{displaymath}
\cbend
(we only define the forward and backward maps on intervals, their
behaviour on $\bot$ is determined.)
\Added{For $\mathrm{add}$, the forward map is the join-preserving conjugate of the meet-preserving map shown
in \exref{intervals-and-maxima-elements}.}
Scaling by negative numbers is
also possible with swapping of bounds, as is multiplication. We will
see an example of the use of these operations in
\secref{language:examples}.

\subsection{Summary}
\label{sec:models:summary}

We have seen that the category $\Fam(\Replaced{\LatConj}{\LatGal})$ has enough structure to
express useful approximation maps at first-order and that
$\Fam(\Replaced{\JoinSLat}{\MeetSLat} \times \JoinSLat^\op)$, which is a Cartesian closed
category with all coproducts, is enough to interpret the total
higher-order language we define in the next section with primitive
types and operations defined in $\Fam(\Replaced{\LatConj}{\LatGal})$. However, we are not
guaranteed by construction that at first-order type, the
interpretations are in fact \Replaced{conjugate pairs}{Galois connections}. We will rectify this
in \secref{definability} using a logical relations construction.

As we did in \secref{approx-as-tangents}, we end the section with a
conjecture relating our categories to Tangent categories.

\Added{\todo{Revisit this conjecture.}}

\begin{conjecture}
  \label{con:tangent-fam}
  (1) The category $\Fam(\MeetSLat)$ is a Tangent category, with the
  tangent bundle
  $T(X, \partial X) = (\Sigma_{x : X}. \partial X(x), \lambda (x,
  \delta x).\, \downset{\delta x})$. (2) The category $\Fam(\Replaced{\LatConj}{\LatGal})$
  is a reverse Tangent category with the analogous definition of
  tangent bundle.
\end{conjecture}

\Replaced{These categories separate points from tangents, somewhat analogous to the situation with
manifolds.}{Comparing this conjecture to \conref{tangent-stable-fns}, we can see
that the difference between $\CM$, $\Lposet$ and $\Fam(\MeetSLat)$,
$\Fam(\LatGal)$ is that the latter have a separation of points from
tangents, somewhat analogous to the situation with manifolds.} If this
conjecture holds, then contrary to the $\Fam(\FinVect)$ representation
of manifolds and differentiable maps, we do not throw away information
about higher derivatives. It is retained in the order structure of the
tangent fibres.

\section{Higher-Order Language}
\label{sec:language}

To demonstrate semiring provenance tracking via automatic differentiation for higher-order programs, we define a simple total functional
programming language, extending the simply-typed lambda calculus. The language is parameterised by a signature
$\Sigma = (\PrimTy, \PrimOp)$ consisting of a set $\PrimTy$ of base types $\rho$ and a family of sets
$\PrimOp^\rho_{\rho_1,\ldots,\rho_n}$ of primitive operations $\phi$ of arity $n$ over those base types.

\subsection{Syntax}
\label{sec:language:syntax}

The syntax is defined in \figref{syntax}. Types includes base types $\rho$ drawn from $\PrimTy$, along with
standard type formers for sums, products, functions and lists% , plus a type constructor $\tyLift$ to allow
% approximation points to be added explicitly, as discussed in \secref{models-of-total-gps}
. Terms include
variables, the usual introduction and elimination forms% , a monadic return and bind for lifted terms
, and
primitive operations $\phi$.

The language is intentionally minimal: it excludes general recursion, and general inductive or coinductive
types, which we will consider in future work (\secref{conclusion}). Typing judgments for terms are standard
and shown in \figref{typing}, with the usual rules for products, sums, functions, and lists.

\input{fig/syntax}
\input{fig/typing}

\subsection{Semantics}
\label{sec:language:semantics}

\input{fig/semantics}

An interpretation of a signature $\Sigma = (\PrimTy, \PrimOp)$ can be given in any category $\cat{C}$ with
finite products, and assigns to each base type $\rho \in \PrimTy$ an object $\sem{\rho}_{\PrimTy}$ in
$\cat{C}$, and to each primitive operation $\phi \in \PrimOp^\rho_{\rho_1,\ldots,\rho_n}$, a morphism
$\sem{\phi}_{\Op}: \sem{\rho_1}_{\PrimTy} \times \ldots \times \sem{\rho_n}_{\PrimTy} \to
\sem{\rho}_{\PrimTy}$.

Assuming that $\cat{C}$ is bicartesian closed and has a list object (\eqnref{lists}), then we can extend an
interpretation of a signature $\Sigma$ to an interpretation of the whole language over
$\Sigma$. \figref{semantics:types} and \figref{semantics:contexts} define the interpretation of types and
contexts as objects of $\cat{C}$ respectively. Terms are interpreted as morphisms between the interpretations
of the context and type, as defined in \figref{semantics:terms}. We have used the notations $\pi_i$ for
projections, $\prodM{f}{g}$ for pairing, $\coprodM{f}{g}$ for (parameterised) copairing, $!_X$ for morphisms
to the terminal object, and $\lambda$ and $\eval$ for currying and evaluation for exponentials.

We have another interpretation $\sem{-}_{\mathit{fo}}$ of the
first-order types (those constructed from primitive types, sums, unit
and products) in any bicartesian category. Such interpretations are
preserved by finite coproduct and coproduct preserving functors:

\begin{lemma}\label{lem:first-order-agreement-types}
  If $\cat{C}$ and $\cat{D}$ are bicartesian and bicartesian closed categories with interpretations of the
  signature $\Sigma$, $F : \cat{C} \to \cat{D}$ is a bicartesian functor, and
  $F(\sem{\rho}_{\PrimTy}) \cong \sem{\rho}_{\PrimTy}$ for all $\rho$, then for all first-order types $\tau$,
  $F(\sem{\tau}_{\mathit{fo}}) \cong \sem{\tau}$, and similar for contexts only containing first-order types.
\end{lemma}

\subsubsection{Interpretation in the Semimodule Model}
\label{sec:language:gps-interpretation}

Given the above, we can now interpret the language in any of the bicartesian closed categories with list
objects we constructed in \secref{models-of-total-gps}. Specifically, we assume that we have an interpretation
of our chosen signature in $\Fam(\SDSemiMod(S))$. Each base type is
assigned a self-dual approximation object, and the first-order type formers preserve self-duality, so every
first-order type denotes a self-dual semimodule. Signatures are
first-order, so it does not matter that $\Fam(\SDSemiMod(S))$ is not closed. Any
such interpretation can be transported to $\Fam(\SemiMod(S))$ along the functor $H$ from \propref{ho-embedding} because it preserves finite products. We
can then interpret the whole language in $\Fam(\SemiMod(S))$.

Interpreting a program $\Gamma \vdash t : \tau$ yields a morphism of $\Fam(\SemiMod(S))$, the
underlying function together with its forward derivative. In contrast to many other bidirectional program
analysis techniques, a single morphism is enough, because a linear map has a transpose for free. At first-order types, \lemref{first-order-agreement-types} identifies the
interpretation with that of the first-order model $\Fam(\SDSemiMod(S))$, whose objects are self-dual, so the
forward derivative is a morphism of $\Fam(\SDSemiMod(S))$ whose conjugate $\conj{f}$ is the backward map
between the same semimodules. 

% Given a model of $\Sigma$ in $\Fam(\LatGal)$, the category $\Fam(\MeetSLat \times \JoinSLat^\op)$ can
% implement $\Sem$. It is bicartesian closed by \corref{mslat-jslat-bcc}, and can model $\tyList$ using infinite
% coproducts (\secref{models-of-total-gps:coproducts-and-products}). The embedding $\HoEmbed: \Fam(\LatGal) \to
% \Fam(\MeetSLat \times \JoinSLat^\op)$ from \propref{ho-embedding} preserves finite products, and so can
% transport each $\sem{\rho}_{\PrimTy}$ and $\sem{\phi}_{\Op}$ into the higher-order setting to provide a model
% of $\Sigma$. Finally the lifting monad in $\Fam(\LatGal)$ (\secref{first-order:lifting}) is also preserved by
% $H$. \roly{don't think we actually show this}

% The correctness of this interpretation is considered in the next section.


\section[Examples]{\added{Examples}}
\label{sec:examples}

\Added{This section runs the interpretation at particular semirings, on particular programs, to show the
analyses that result. Each example is mechanised in the accompanying Agda development.}

\subsection[Perturbation Bounds]{\added{Perturbation Bounds}}
\label{sec:examples:perturbation-bounds}

\Added{The first two instances demonstrate the perturbation-bound reading of
\secref{approx-as-tangents:richer-semirings}: the analysis reports how much each output can change, rather
than whether it can. Sections \ref{sec:examples:absolute-bounds} and \ref{sec:examples:relative-bounds}
give the two presentations, absolute and relative, run on the database $\mathit{db} = [(\mathsf{a}, 3),
(\mathsf{b}, 1), (\mathsf{a}, -3)]$ of \secref{introduction:galois-slicing}. In both, the approximation
object for a number is the free semimodule $S^2$, one bound for each direction of change.}

\subsubsection[Absolute bounds]{\added{Absolute bounds}}
\label{sec:examples:absolute-bounds}

\Added{At the min-plus semiring, a pair of scalars $(d, u)$ records that a value may fall by at most $d$ and
rise by at most $u$, so that a number $x$ is known to lie in the interval $[x - d, x + u]$; the pair
$(\infty, \infty)$ carries no information.}

\Added{Multiplication's derivative carries no information at this semiring, since no translation bounds a
scaling (\secref{approx-as-tangents:richer-semirings}), so we run the additive part of the running example:
the selection-and-sum query $\mathrm{query}\,l\,\mathit{db} = \mathrm{sum}\,[ q \mid (l', q) \leftarrow
\mathit{db}, l \equiv l' ]$, with $\mathrm{query}\,\mathsf{a}\,\mathit{db}
= 0$. Addition's derivative entries are translations by $0$, and selection's entries are the semiring's $0$
and $1$, as before. The backward derivative at this run is a linear map}
\begin{AddedBlock}
\begin{displaymath}
  \partial(\mathrm{query}\,l\,\mathit{db})_r : S^2 \multimap S^2 \times S^2 \times S^2
\end{displaymath}
\end{AddedBlock}
\noindent\Added{with one factor for each number in the database (the frozen label positions and the list terminator
contribute trivial factors, omitted here). Because the input is a list, the type of the derivative depends on
the shape of the input.}

\Added{Backwards, asking for the output to lie within $[-\frac{1}{10}, \frac{1}{10}]$, that is, bounds
$(\frac{1}{10}, \frac{1}{10})$, returns the same bounds at both $\mathsf{a}$-numbers and $(\infty,
\infty)$ at the $\mathsf{b}$-number: changing either $\mathsf{a}$-number alone by at most $\frac{1}{10}$
keeps the output within its interval, while the $\mathsf{b}$-number is unconstrained. Forwards, if the first
$\mathsf{a}$-number is known to within bounds $(\frac{1}{2}, 0)$ and the second to within $(\frac{1}{5},
\frac{1}{2})$, the output is known to within $(\frac{1}{5}, 0)$: bounds for the same position combine by
$\min$, reflecting the reading in which one input changes at a time.}

\subsubsection[Relative bounds]{\added{Relative bounds}}
\label{sec:examples:relative-bounds}

\Added{In the second presentation a scalar bounds the relative rather than the absolute change of a value,
and the two readings trade places over the arithmetic operations: relative bounds pass through multiplication
unchanged, but are amplified by addition. Because the absolute presentation could say nothing about
multiplication, it exercised only the additive part of the running example. With relative bounds we can run
the full weighted-sum query $\mathrm{total}$ of \secref{introduction:galois-slicing}, prices included, on
the same database.}

\Added{We instantiate at the min-times semiring. The derivative of multiplication now has entries $1$ on both
arguments, since a bound on the relative error of either factor is a bound on the relative error of the
product. The derivative of addition at values $(x, y)$ has entry $|x| / |x + y|$ on its first argument and
symmetrically on its second: the factor by which the relative error of a summand is amplified in the sum. When
$x + y = 0$ the factor is infinite, and the entry is the semiring zero $\infty$.}

\Added{In the run of $\mathrm{total}\,\mathsf{b}$, multiplication behaves well: if the price of
$\mathsf{b}$ is known to within $10\%$, the output is too, and likewise for the selected quantity;
backwards, an output bound of $10\%$ requires the used inputs to within $10\%$ and constrains nothing else.
The run of $\mathrm{total}\,\mathsf{a}$ meets the cancellation: its two contributions, $2 \cdot 3$ and $2
\cdot (-3)$, sum to zero, the amplification factor of the final addition is infinite, and no relative bound
on either $\mathsf{a}$-quantity, however tight, yields a bound on the output. The cancellation that the
Boolean instance reported as a spurious dependency in \secref{introduction:galois-slicing} is here reported
as an infinite condition number, in the sense of numerical analysis~\cite{high:ASNA2}.}

\Added{Neither reading of perturbation bounds is better than the other. Absolute bounds propagate
accurately through addition but carry no information through multiplication, while relative bounds propagate
accurately through multiplication but are amplified by addition, and lost entirely where contributions
cancel. The choice between the two presentations depends on what the analysis is meant to observe.}

\subsection[Linked Inputs and Outputs]{\added{Linked Inputs and Outputs}}
\label{sec:examples:linked}

\Added{When outputs overlap in their dependencies, the two derivatives compose into a further analysis: a
map sending a selection of outputs to the outputs related to it by a shared dependency. The
\emph{moving average} is the standard example of such overlap. With window two and four inputs,}
\begin{AddedBlock}
\begin{displaymath}
  \mathrm{mavg}\,(x_1, x_2, x_3, x_4) =
  (\tfrac{1}{2} (x_1 + x_2),\; \tfrac{1}{2} (x_2 + x_3),\; \tfrac{1}{2} (x_3 + x_4))
\end{displaymath}
\end{AddedBlock}
\noindent\Added{Each adjacent pair of outputs shares an input, and the
first and last outputs share none. Running over the Booleans at the input $(1, 2, 4, 8)$, the Jacobian and
the composite of the backward and forward derivatives are}
\begin{AddedBlock}
\begin{displaymath}
  \partial_{\Two}(\mathrm{mavg}) =
  \begin{pmatrix} 1 & 1 & 0 & 0 \\ 0 & 1 & 1 & 0 \\ 0 & 0 & 1 & 1 \end{pmatrix}
  \qquad
  \partial_{\Two}(\mathrm{mavg}) \comp \transpose{\partial_{\Two}(\mathrm{mavg})} =
  \begin{pmatrix} 1 & 1 & 0 \\ 1 & 1 & 1 \\ 0 & 1 & 1 \end{pmatrix}
\end{displaymath}
\end{AddedBlock}
\Added{The composite relates each output to the outputs it shares an input with: selecting the first output
returns the first and second, and the third stays $\bot$ because the two windows are disjoint. This is the
\emph{cognacy} analysis of linked visualisations~\cite{perera22,bond25,psallidas18}, computed here as a
composite of the two derivatives of one run.}

\subsection[Counting]{\added{Counting}}
\label{sec:examples:counting}

\Added{The counting semiring $(\mathbbm{N}, +, 0, \cdot, 1)$ of
\secref{approx-as-tangents:richer-semirings} refines dependency tracking from whether an input is used to
how many times. Every argument of every operation is given the coefficient $1$, so a Jacobian entry counts
the paths from an input position to an output position.}

\Added{On the weighted-sum query of \secref{introduction:galois-slicing}, the backward derivative of the
output returns a count for every position: $1$ at each selected quantity, $0$ at the unselected row and the
other price, and $2$ at the price of the queried label, which is used once in each selected row. The count
of $2$ survives the cancellation that zeroes the rational derivative, because the counting semiring is
positive and counts cannot cancel. The three instances thus answer three different questions at the same
position: the rational Jacobian reports sensitivity ($0$), the Boolean Jacobian reports possible dependence
($1$), and the counting Jacobian reports usage ($2$). Zero-testing sends the count $2$ to the Boolean $1$
exactly, as \secref{approx-as-tangents:richer-semirings} leads us to expect of a positive semiring, whereas
the same specialisation applied to the rational entry loses the dependency.}

\section{Correctness of the Higher-Order Interpretation}
\label{sec:definability}

The interpretation of the higher-order language at higher types mixes
the interpretation of the underlying computation with the
interpretation of the semiring dependency matrices
(c.f. \thmref{fam-closed}). It is therefore not immediate that
intepreting a program in $\Fam(\SemiMod(S))$ will give us the expected
interpretation purely on first order data. Note that the fullness of
the functor $H : \Fam(\SDSemiMod(S)) \to \Fam(\SemiMod(S))$ does mean
that we are guaranteed that the tangent maps are correctly computed.
\citet{vakar22} and \citet{nunes2023} construct custom
instances of categorical sconing arguments to prove correctness of
their higher-order interpretation with respect to normal
differentiation. Instead of doing this, we make use of a general
\emph{syntax free} theorem due to \citet{fiore-simpson99}. The proof
of this depends on the construction of a Grothendieck Logical Relation
over the extensive topology on the category $\cat{C}$, but the
statement of the theorem does not rely on this. We have formalised
this proof in Agda (see \texttt{conservativity.agda} in the
supplementary material).

\begin{theorem}[\citet{fiore-simpson99}]
  \label{thm:glr-definability}
  Let $\cat{C}$ be an extensive bicartesian category, $\cat{D}$ be a
  bicartesian closed category, and $F : \cat{C} \to \cat{D}$ a functor
  preserving finite products and coproducts. Then there is a category
  $\GLR(F)$ and functors $p : \GLR(F) \to \cat{D}$ and
  $\hat{F} : \cat{C} \to \GLR(F)$, such that:
  \begin{enumerate}
  \item $\GLR(\cat{D},F)$ is bicartesian closed;
  \item $F = p \circ \hat{F} : \cat{C} \to \cat{D}$;
  \item The functor $p$ strictly preserves the bicartesian closed structure; and
  \item The functor $\hat{F}$ is full and preserves the bicartesian structure.
  \end{enumerate}
\end{theorem}

\begin{remark}
  Compared to the exact result stated at the end of
  \citet{fiore-simpson99}'s paper, we have made two modifications,
  justified by our Agda proof. First, we generalise to the case where
  $\cat{C}$ is not Cartesian closed, and the functor $F$ does not
  preserve exponentials. Examination of the proof reveals that if this
  is the case, then $\hat{F}$ also preserves exponentials, but it is
  not needed for the result stated. Second, Fiore and Simpson restrict
  to the case when $\cat{C}$ is small to be able to construct
  Grothendieck sheaves on this category. We use Agda's universe
  hierarchy to simply construct ``large'' sheaves at the the
  appropriate universe level.
\end{remark}

We use \thmref{glr-definability} to show that the
interpretation of the language in the category $\Set$ agrees with the
higher-order interpretation in $\Fam(\SemiMod(S))$
on the underlying function at first order, as required. This shows that the
higher-order interpretation does what we expect in the underlying
interpretation of terms, and that the approximation information does
not interfere.

\begin{theorem}
  \label{thm:underlying-interp-equal} For all
  $\Gamma \vdash M : \tau$, where $\Gamma$ and $\tau$ are first-order,
  the underlying function in the interpretation
  $\sem{\Gamma \vdash M : \tau}_{\Fam(\SemiMod(S))}$ is equal to the
  interpretation $\sem{\Gamma \vdash M : \tau}_{\Set}$ in $\Set$.
\end{theorem}

\begin{proof}
  Instantiate \thmref{glr-definability} with the functor
  \begin{displaymath}
    \langle \mathrm{Id} , \pi_1 \rangle : \Fam(\SemiMod(S)) \to \Fam(\SemiMod(S)) \times \Set
  \end{displaymath}
  that is the identity in the first component and projects out the
  underlying function in the second. By
  \lemref{pi1-preserve-products-and-coproducts}, this functor
  preserves products and coproducts. For each
  $\Gamma \vdash M : \tau$, we obtain a $g$ such that
  $g = \sem{\Gamma \vdash M : \tau}_{\Fam(\SemiMod(S))}$ and
  $\pi_1(g) = \sem{\Gamma \vdash M : \tau}_\Set$. Substituting $g$
  yields the result.
\end{proof}


\section{Related Work}
\label{sec:related-work}

\begin{DeletedBlock}
\paragraph{Stable Domain Theory}

\roly{Now covered in Future Work.} Stable Domain Theory was originally proposed by \citet{berry79} as a
refinement of domain theory aimed at capturing the intensional behaviour of sequential programs, and
elaborated on subsequently by \citet{berry82} and \citet{amadio-curien}. Standard domain-theoretic models
interpret programs as continuous functions, preserving directed joins; Berry observed that this continuity
condition alone is too permissive to model sequentiality. Stability imposes additional constraints to reflect
how functions preserve bounded meets of approximants, effectively requiring that the evaluation of a function
respect a specific computational order. Though stable functions do not fully characterise sequentiality,
because they admit $\mathrm{gustave}$-style counterexamples (\exref{parallel-or}), they remain an appropriate
notion for studying the sensitivity of a program to partial data at a specific point.

Our use of Stable Domain Theory diverges from the traditional aim of modelling infinite or partial data,
however. Instead, we follow a line of work that uses partiality as a qualitative notion of approximation
suitable for provenance and program slicing (discussed in more detail in \secref{related-work:galois-slicing}
below). Paul Taylor’s characterisation of stable functions via local Galois connections on principle downsets
provides the semantic underpinning for the reverse maps used in Galois slicing~\cite{taylor99}. Our work
builds on these ideas by interpreting Galois slicing as a form of differentiable programming, using the
machinery of CHAD to present Galois slicing in a denotational style.
\end{DeletedBlock}

\paragraph{Automatic Differentiation}

Automatic differentiation (AD), discussed in \secref{approx-as-tangents:autodiff}, is the idea of computing
derivatives of functions expressed as programs by systematically applying the chain rule. The observation that
these derivative computations could be interleaved with the evaluation of the original program is due to
\citet{linnainmaa76}, who showed how the forward derivative $\pushf{f}_x$ of $f$ at a point $x$ could be
computed alongside $f(x)$ in a single pass, dramatically improving the efficiency of derivative evaluation
over symbolic or numerical differentiation. This insight became the foundation of forward-mode AD, which
underpins many optimisation and scientific computing tools, including JAX~\cite{jax2018github}.
\citet{griewank89} showed how the Wengert list, the linear record of assignments used in forward-mode to
compute derivatives efficiently, could be traversed in reverse to compute the pullback map. This two-pass
approach is the foundation of reverse-mode AD, and closely resembles implementations of Galois slicing
(\secref{related-work:galois-slicing} below) that record a trace during forward slicing for use in backward
slicing.

More recent approaches to automatic differentiation have emphasised semantic foundations. \citet{elliott18}
proposed a categorical model of AD that interprets programs as functions enriched with their derivatives,
giving a compositional account of differentiation based on duality and linear maps. Vákár and
collaborators~\cite{vakar22,nunes2023} developed the CHAD framework which inspired this paper, using
Grothendieck constructions over indexed categories to capture both values and their tangents in a
compositional semantic structure. These perspectives shed light on the categorical structure of AD and guide
the design of systems that generalise AD% beyond real analysis
, including the application to data provenance
and slicing explored in this paper.

\paragraph{\GPS}
\label{sec:related-work:galois-slicing}

Galois slicing was introduced by Perera and collaborators~\cite{perera12a,perera13} as an operational approach
to program slicing for pure functional programs, based on Galois connections between lattices of input and
output approximations. \Deleted{A connection to stable functions in relation to minimal slices for
short-circuiting operations was alluded to in~\citet{perera13}, but not explored.} Subsequent work extended
the approach to languages with assignment and exceptions~\cite{ricciotti17} and
\Replaced{concurrency}{concurrent systems, applying Galois slicing to the $\pi$-calculus}~\cite{perera16d}.
\Deleted{For the $\pi$-calculus the analysis shifted from functions to transition relations, considering
individual transitions $P \longrightarrow Q$ between configurations $P$ and $Q$ as analogous to the edge
between $x$ and $f(x)$ in the graph of $f$, and building Galois connections between $\downset{P}$ and
$\downset{Q}$.} \Replaced{The present work is instead denotational, recasting dependency analysis as
differentiation over a semiring; the Boolean case recovers Galois slicing, and other semirings give further
analyses. Galois slicing is also more general than our framework in two respects that we plan to consider in
future work (\secref{conclusion}): it computes \emph{program slices} from approximation lattices representing
partially erased programs, and it approximates the \emph{shape} of data rather than only the values at fixed
positions.}{The main difference with the approach presented here is that the earlier work also computes
\emph{program slices}, using approximation lattices that represent partially erased programs; we discuss this
further in \secref{conclusion} below.} More recently, \citet{perera22} presented a variant of Galois slicing
over Boolean algebras\Replaced{, which unlike plain lattices are complemented, giving each analysis a
conjugate}{ rather than plain lattices}. \Replaced{Composing an analysis with its conjugate computes
\emph{related outputs}, which they used to support interactive visualisations with linked selections.
\citet{bond25} revisited this using \emph{dynamic dependence graphs}, computing the conjugate from the
opposite graph.}{In this setting every Galois connection $f \dashv g: A \to B$ has a conjugate $g^\circ \dashv
f^\circ: B \to A$, where $f^\circ$ denotes the De Morgan dual $\neg \comp f \comp \neg$~\cite{jonsson51}. The
provenance analysis can then be composed with its own conjugate to obtain a Galois connection which computes
\emph{related outputs} (e.g., selecting a region of a chart and observing the regions of other charts which
share data dependencies). \citet{bond25} revisited this approach using \emph{dynamic dependence graphs} to
decouple the derivation of dependency information from the analyses that make use of it, and observing that to
compute the conjugate analysis one can just use the opposite graph.} \Added{The related-outputs analysis is
the cognacy relation shown in \secref{examples:linked}, computed here by composing a Boolean Jacobian with its
transpose.}

\paragraph{\added{Semiring Provenance}}

\Added{Provenance for database queries has its own semiring tradition, originating with the \emph{provenance
semirings} of \citet{green07}: a query result carries an annotation from a commutative semiring, and the free
semiring $\mathbbm{N}[X]$ of provenance polynomials is universal, specialising to other semirings by
evaluation. Our framework draws on the same semirings but computes a \emph{derivative} rather than an
annotation: a Jacobian over the semiring relating input positions to output positions, in place of a value
attached to the output. A derivative is a linear map, so the provenance it records is \emph{linearised}
(\secref{approx-as-tangents:linearised-provenance}), and joint use appears as a coefficient rather than the
degree of a polynomial. Over a non-positive semiring, the Boolean dependency our framework reports is a sound
over-approximation, since the exact analysis can cancel a contribution to zero where the Boolean composite
still records a dependency. Extending provenance to deletion and difference likewise calls for non-positive
annotations~\cite{green09,amsterdamer11,geerts10}.}

\paragraph{\added{Discrete and Qualitative Derivatives}}

\Added{Discrete data has its own differential calculi, developed independently of numerical differentiation. The
\emph{Boolean differential calculus} of \citet{steinbach17} arose in the 1950s from the detection of faults in
digital circuits; its \emph{simple derivative} $\partial f / \partial x_i = f|_{x_i = 0} \veebar f|_{x_i = 1}$,
the exclusive-or of $f$ with the input $x_i$ set to $0$ and to $1$, is $1$ exactly when flipping that input
flips the output. Exclusive-or and conjunction make the Booleans a \emph{ring} rather than a lattice, so this is
a derivative over the two-element field, the instance $S = \mathbbm{F}_2$ of our construction. It differs from
the Boolean case of Galois slicing (\secref{related-work:galois-slicing}) only in the semiring. Because
exclusive-or cancels, the derivative is exact, whereas the join of the Boolean lattice only reports a possible
dependency. The sign domain gives a second such calculus. The rule of signs of abstract
interpretation~\cite{cousot77} and the \emph{qualitative derivatives} of qualitative
physics~\cite{dekleer84,kuipers86} ask whether increasing an input increases, decreases, or cannot affect an
output, and are the sign semiring of \secref{approx-as-tangents:qualitative}. Our account treats the semiring as
a parameter, placing these calculi alongside automatic differentiation and provenance as instances of a single
construction.}

\paragraph{Tangent Categories and Differential Linear Logic}

\emph{Tangent categories}, due originally to \citet{rosický84} and developed by \citet{cockett14,cockett18},
provide an abstract categorical framework for reasoning about differentiation, inspired by the structure of
the tangent bundle in differential geometry. In a tangent category, each object $X$ is equipped with a tangent
bundle $T(X)$, and each morphism $f: X \to Y$ has a corresponding differential map $T(f): T(X) \to T(Y)$
satisfying axioms analogous to the chain rule and linearity of differentiation. Tangent categories generalise
Cartesian differential categories \cite{cdcs}, which model differentiation over Cartesian closed categories
using a syntactic derivative operator. Reverse Tangent categories \cite{reverse-tangents} further axiomatise
the existence of reverse derivatives. \Replaced{Our construction realises differentiation of this kind concretely, over an arbitrary commutative
semiring, and we conjecture that the categories underlying it are examples of tangent, or reverse tangent,
categories.}{In \conref{tangent-stable-fns} and
\conref{tangent-fam}, we have conjectured that the categories we have identified in this paper as models of
\GPS are Tangent categories. This would clarify the role of higher derivatives in \GPS, which we conjecture
are related to \emph{program differencing}.} There are likely links to Differential Linear Logic
\cite{ehrhard_differential_2006}. Differential Linear Logic and the Dialectica translation have been used to
model reverse differentiation by \citet{kerjean-pedrot2024}.


% \paragraph{Lens Categories}

% \cite{spivak19}

\section{Conclusion and Future Work}
\label{sec:conclusion}

We have presented a semantic account of data provenance as a form of automatic differentiation, representing
the dependency of outputs on inputs as a derivative whose Jacobian is a matrix over a commutative semiring.
Traditional automatic differentiation, and various program analysis techniques such as \GPS, arise as
instances. The model interprets an expressive higher-order language for querying and manipulating data, and
reveals new applications such as approximation by intervals~(\secref{approx-as-tangents:perturbation-bounds}).
It also opens up an opportunity to integrate traditional symbolic data provenance techniques with the AD-based
methods used in explainable AI, such as the saliency maps of Grad-CAM~\cite{selvaraju20}. Our categorical
approach admits a modular construction of the model, and the use of general theorems, such as Fiore and
Simpson's \thmref{glr-definability}, to prove properties of the interpretation. We have focused on
constructions that enable an executable implementation in Agda.

\paragraph{Alternative semiring models}

The choice of scalar semiring is the principal degree of freedom in our framework, and choices
beyond those studied here seem worth pursuing. The perturbation-bound semirings of
\secref{examples:perturbation-bounds} point towards metric approximation: a Lawvere metric space is a
category enriched over the min-plus quantale, whose operations are those of our absolute-perturbation
semiring, so metric spaces provide a native notion of approximation already partly within reach of the
framework. Relating this to \citet{edalat-heckmann98}'s computational model of metric spaces is one promising
direction. A second is to replace the setoids carrying values with sets equipped with a semiring-valued
equality relation, over a partially ordered semiring with a residual for its multiplication.

% \paragraph{Categorical Models of Differentation} We have already identified \conref{tangent-stable-fns} and
% \conref{tangent-fam} as potential connections to categorical models of differentation.

\paragraph{Connection to sequentiality and stable functions}

In our model the forward derivative carried by a morphism of $\Fam(\SemiMod(S))$ is supplied as part
of the interpretation, not determined by the underlying function, so it need not reflect how that function
depends on its input. Berry's
\emph{stable functions}~\cite{berry79,berry82}, by contrast, come equipped with an intrinsic forward and
backward derivative. Stability was proposed as a refinement of domain theory aimed at capturing the
intensional behaviour of sequential programs, imposing constraints on how functions preserve bounded meets of
approximants; although stable functions do not fully characterise sequentiality, because they admit
$\mathrm{gustave}$-style counterexamples, they remain an appropriate notion for studying the sensitivity of a
program to partial data at a specific point~\cite{amadio-curien}. Paul Taylor's characterisation of stable
functions via local Galois connections on principal downsets provides a semantic underpinning for the reverse
maps of Galois slicing~\cite{taylor99}. The stable-function story has a tight connection to Galois slicing but a looser one to automatic
differentiation, so we intend to develop it in separate work, relating Galois slicing to stable functions
rather than to differentiation. Such a development would also clarify the links to Differential Linear
Logic~\cite{ehrhard_differential_2006} and to Crubill\'e's identification of stable functions with power
series~\cite{crubille18}, relating the qualitative derivatives studied here to the analytic derivatives of
that setting.

\paragraph{Shape Approximation}

Lists are the only recursive datatype in our source language, so supporting general inductive and
coinductive types is important future work, perhaps adapting \posscite{nunes2023} automatic differentiation
for $\mu\nu$-polynomial functors. Richer datatypes bring the question of \emph{shape approximation}:
approximating not only the values at the positions of a structure but the structure itself, such as the
length of a list or the depth of a tree. Shape approximation features in earlier work on Galois
slicing~\cite{perera12a,perera16d,ricciotti17}, whose approximation lattices represent partially erased data,
whereas our model holds the shape of each input fixed. Our Jacobians are spans, and polynomial functors are
the shape-dependent generalisation of spans~\cite{gambino-kock}; the finite polynomials form the Lawvere
theory for commutative semirings, the algebra of our scalars. The corresponding datatypes are containers,
whose derivative is the type of one-hole contexts, carrying a notion of linear map and a chain rule for both
inductive and coinductive types~\cite{abbott05}. This already gives a set-theoretic account of shape
approximation; weighting positions with our semiring scalars would extend our analyses to varying shapes.

\paragraph{Recursion and Partiality}  This work treats only total programs, and makes no use of directed completeness or similar
properties of domains. Extending to recursion and partiality would introduce a notion of undefinedness, and
hence a definedness order to be reconciled with any approximation order on the same data. Non-termination and
infinite data are both partiality phenomena, approximating unbounded computation by finite observation, so
potentially the finite prefixes of an infinite or non-terminating computation would fall under the same
extension.

\paragraph{Source-To-Source Translation Techniques}

% To formalise this, we'd need to... See also $\Fam(\PSh_{\CMon}(\namedcat{Syn}))$?

An interesting alternative to the denotational approach presented here, and to the trace-based approaches used
in some implementations of provenance, would be to develop a
source-to-source transformation, in direct analogy with
the CHAD approach to automatic differentiation \cite{vakar22,nunes2023}. In their approach, forward and
reverse-mode AD are implemented as compositional transformations on source code, guided by a universal
property: they arise as the unique structure-preserving functors from the source language to a suitably
structured target language formalised as a Grothendieck construction. Adapting this to our setting
would allow the analysis to be ``compiled in'', avoiding the need for
a custom interpreter and potentially exposing opportunities for optimisation.


\bibliographystyle{ACM-Reference-Format}
\bibliography{bib}

% \pagebreak
% \appendix
% \input{appendix/notes}

\end{document}
